\chapter{Anatomie und Physiologie}

\emph{Maintainer: Sebastian Gabriel}


\gAuthNotes{
Changelog: 

\begin{description}
    \item[2009-02-25] Start Changelog.
    \item[2009-05-06] GABS: Korrekturen von KOCH
    eingearbeitet.
\end{description}
}

\minitoc

\section{Die Zelle}

Die Zelle ist die kleinste lebende Einheit in unserem
Körper. Sie besteht u.a. aus einem Zellkern, der die
gesammten Erbinformationen enthält und verschiedenen
Zellorganellen. Sie wird von der Zellmembran umhüllt. 

Artverwandte Zellen bilden im Körper meist einen
Zellverband, welcher \gT{Gewebe}\index{Gewebe} genannt wird. 

\begin{figure}[H]
   
   
   
   
   
    \includegraphics[width=\linewidth]{./images/anatomie/Animal_cell_structure_de-edited.png}
    \captionof{figure}{\gAuthNotesPicLegalOk{}{PD}Eine
    Zelle, Schema. \gSourcePicture{Mariana Ruiz Villarreal,
    Public domain}}
\end{figure}



\section{Das Skelett und der Bewegungsapparat}
\begin{figure}[H]
    \includegraphics[width=10cm]{./images/anatomie/Human_skeleton_front_de.png} \captionof{figure}{\gAuthNotesPicLegalOk{}{ PD }Das Menschliche Skelett.
    \gSourcePicture{Mariana Ruiz Villarreal, Public domain}}
    % TODO Pic Skelett Fußwurzelknochen eindeutsche
\end{figure}

Das Skelettsystem besteht aus Knochen und dient 
haupts\"achlich als

\begin{itemize}
    \item St\"utzger\"ust  des K\"orpers
    \item Ursprungs-  und Ansatzfl\"ache f\"ur die Muskulatur
\end{itemize}

und weiters als

\begin{itemize}
    \item Kalziumspeicher ($Ca^{++}$)
    \item Blutbildungsorgan (rotes Knochenmark!)
\end{itemize}

\subsection{Der Knochen}

Je nach ihrer Form unterteilt man die Knochen in: 
\begin{description}
    \item[\gT{Platte Knochen}] (Beckenknochen,
    Schulterblatt),
    \item[\gT{Röhrenknochen}] (Oberarmknochen, Schienbein)
    und
    \item[\gT{Würfelknochen}] (Hand-,
    Fußwurzelwurzelknochen)
\end{description}

\paragraph{Aufbau eines Knochens}~

\gAuthNotes{GABS \& KOCH: evt. umstellen auf deutsch -- latein, aber was heisst
Kortikalis auf deutsch?}

\gMarried{

\begin{description}
    \item[\gT{Periost}] \index{Periost}\gE{Beinhaut}, \"uberzieht den Knochen und ist sehr gut mit
Nerven versorgt.  Verantwortlich f\"ur Dickenwachstum und Frakturheilung
\item[\gT{Kortikalis}] \index{Kortikalis}feste \"au{\ss}ere
Schicht, \"au{\ss}ere Form
\item[\gT{Spongiosa}] \index{Spongiosa}\index{Knochenb\"alkchen}Knochenb\"alkchen ${\rightarrow}$ Zug- und Druckfestigkeit
\item[\gT{Markhöhle}] s.u.
\item[\gT{Epiphysenfuge}] Wachstumsfuge
\item[\gT{Gelenkfläche}] mit knorpeligen Überzug als Teil
eines Gelenkes. Manche Knochen haben keine Gelenksfläche.
\end{description}

Lange Röhrenknochen kann man ausserdem in 
\begin{compactitem}
    \item \gT{Kopf},
    \item \gT{Hals} und
    \item \gT{Schaft}
\end{compactitem}
unterteilen.


}
{\includegraphics[width=\linewidth]{./images/anatomie/roehrenknochen-querschnitt.png}
\captionof{figure}{Ein langer Röhrenknochen im Querschnitt.
\gSourcePicture{Lena Hirtler, modifiziert.}} }

\subsubsection{Das Knochenmark}

\paragraph{Rotes Knochenmark ist f\"ur die Blutbildung
zust\"andig}

Das rote Knochenmark ist das blutbildende Organ
schlechthin. Hier werden \gT{Rote Blutk\"orperchen}
(\gT{Erythrozyzen}), \gT{wei{\ss}e Blutk\"orperchen} (\gT{Leukozyten}) sowie
\gT{Blutpl\"attchen} (\gT{Thrombozyten}) gebildet
(\gSee{topic:blut}). Bei der Geburt findet man es in
s\"amtlichen Knochen, im Laufe des Lebens bildet es sich zur\"uck und verwandelt sich in {\quotedblbase}Gelbes Fettmark{\textquotedblleft}. Rotes Knochenmark findet man dann lediglich in den platten Knochen wie Becken, Brustbein
und Sch\"adel, wo es f\"ur die Blutbildung verantwortlich
ist und normalerweise zeitlebens erhalten bleibt. 

\gAuthNotes{KOCH: wichtig? ev. weglassen?

GABS: Ja, wichtig wegen KT, onko-Patienten, etc. Soll
bleiben. 

KOCH: ok.} 

Einige medizinische Behandlungsmethoden
verursachen als Nebenwirkung eine massive \gEs{Sch\"adigung
des Knochenmarks} (z.B. Bestrahlung, Chemotherapien bei
Krebspatienten ("`onkologische
Patienten"'), \ldots). Diese Patienten leiden neben einer
Blutarmut auch an einem schwachen Immunsystem, da nicht gen\"ugend Wei{\ss}e Blutk\"orperchen gebildet werden. Sie k\"onnen daher sehr
leicht schwerwiegende Infektionen bekommen. Diese Patienten
sind im Krankentransport h\"aufig anzutreffen. 

\paragraph{Gelbes Knochenmark besteht aus Fett} 

Im Laufe des Lebens bildet es sich das rote Knochenmark in den
R\"ohrenknochen zu fettreichem, gelben Knochenmark um. 

Das gelbe Fettmark kann bei schweren Knochenverletzungen von Bedeutung
sein, da verletzungs\-bedingt kleine Fettptr\"opchen in die Blutbahn
gelangen k\"onnen, die dann als Fettgerinnsel  zu ernsten
Komplikationen f\"uhren k\"onnen
(Fettembolie\gAuthNotes{+ KOCH}).

\subsubsection{Knochenwachstum}~

Das Skelett eines Neugeborenen besteht bei der Geburt zum gr\"o{\ss}ten
Teil aus Knorpel. Erst mit zunehmendem Lebensalter beginnt dieser von
Knochenkernen ausgehend, zu verkn\"ochern. Eine
\gEs{Wachstumszone}
(\index{Epiphysenfuge}\gT{Epiphysenfuge}) bleibt dabei allerdings erhalten.

Die Epiphysenfuge enth\"alt beim Kind/Jugendlichen bis
etwa zum Ende der Pubert\"at \gE{teilungsf\"ahiges
Knorpelgewebe}, daher ist ein \gEs{L\"angenwachstum}
m\"oglich. Beim Erwachsenen sind diese knorpeligen Fugen
hingegen vollst\"andig \gE{verkn\"ochert}, soda{\ss}
keinerlei Wachstum mehr stattfinden kann.

\subsubsection{Gelenke}

\gDefinition{(Meistens bewegliche) Verbindung zwischen
mindestens zwei Knochen.}

\paragraph{Prinzipieller Aufbau}~

\gMarried{
\begin{itemize}
\item Ein Gelenk ist von einer bindegewebigen Kapsel umgeben und durch
Sehnen und B\"ander stabil verbunden.

\item Die Gelenksfl\"achen sind zur Reibungs\-ver\-minderung mit 
Knorpel \"uberzogen und der Gelenksspalt mit einer d\"unnen Schicht 
Schmierfl\"ussigkeit \gZ{(Synovialfl\"ussigkeit)} gef\"ullt. 

\end{itemize}}
{~

\includegraphics[width=\linewidth]{./images/anatomie/anatomie-gelenk-edited.png}
\captionof{figure}{\gAuthNotesPicLegalOk{}{Lena }Ein Gelenk
im Querschnitt. \gSourcePicture{Lena
Hirtler}}\gAuthNotes{KOCH. löschen: Membrana synovialis
und membrana fibrosa}}


\paragraph{Es gibt verschiedene Arten von Gelenken}

Klassische Beispiele: \index{Kugelgelenk}\gT{Kugelgelenk} (H\"uftgelenk), \index{Scharniergelenk}\gT{Scharniergelenk}
(Fingergliedgelenke), \index{Sattelgelenk}\gT{Sattelgelenk}
(Daumengrundgelenk).  


\begin{minipage}{\linewidth}
\includegraphics[width=\linewidth]{./images/noauth/AASdev20090228-img23.png}
\captionof{figure}{\gAuthNotesPicLegalNotOk{img23}{Arten von
Gelenken}{Arten von
Gelenken}}  
\end{minipage}


\subsection{Der Sch\"adel} 

Kopf (\gT{Caput}) / Sch\"adel (\gT{Cranium})

 


\begin{gWideParEnv}
\gMarried{
\includegraphics[width=\linewidth]{./images/anatomie/Human_skull_side_simplified_hirngesicht_de-edited.png}
\captionof{figure}{\gAuthNotesPicLegalOk{}{PD}
Der Sch\"adel gliedert sich in den \gT{Hirnsch\"adel} und
den \gT{Gesichtssch\"adel}}. 
\gSourcePicture{Mariana Ruiz Villarreal, Public Domain}
}
{
\includegraphics[width=\linewidth]{./images/anatomie/Human_skull_front_simplified_bones_de-edited.png}
\captionof{figure}{\gAuthNotesPicLegalOk{}{PD}Schädel Vorderansicht.
\gSourcePicture{Mariana Ruiz Villarreal, Public Domain}}
% \includegraphics[width=\linewidth]{./images/anatomie/Human_skull_side_simplified_bones_de-edited.png}
% \captionof{figure}{\gAuthNotesPicLegalOk{}{PD}Schädel seitlich.
% \gSourcePicture{Mariana Ruiz Villarreal, Public Domain}}
}\end{gWideParEnv}

\gA{STEU 2009-06-28: Genaue Bezeichnungen der einzelnen
Schädelknochen wirklich notwendig? Unterkiefer, Oberkiefer,
Kalotte.

EMHO: Schrift grau. }

\gMarried{\begin{itemize}
\item Sch\"adelbasis
\item Sch\"adeldach (Kalotte)
\item Die platten Knochen sind durch Sch\"adeln\"ahte  
verzahnt miteinander verkn\"ochert
\end{itemize}
}
{
~
}



\subparagraph{Fontanelle}~
\index{Fontanelle}

\gMarried{Die \gT{Fontanellen} sind \gEs{knorpelige 
Verbindungen} zwischen Sch\"adelknochen beim Neu\-geborenen
(der Sch\"adel mu{\ss} sich f\"ur Geburt verformen
k\"onnen!). \gZ{Sie schlie{\ss}en  sich bis zum 18.
Lebensmonat,} verkn\"ochern  aber erst vollst\"andig mit
zunehmendem  Lebensalter \gZ{(ca. 20 Lebensjahr
abgeschlossen)}.

\begin{itemize}
\item \gZ{gro{\ss}e Fontanelle (rot): zw. Stirnbein u.
Scheitelbeinen}
\item \gZ{kleine Fontanelle (gr\"un): zw. Scheitelbeinen u.
Hinter\-haupts\-bein}
\end{itemize}
}
{
\includegraphics[width=\linewidth]{./images/anatomie/fontanelle-dt.jpg}
\captionof{figure}{\gAuthNotesPicLegalOk{Fontanelle}{Gray}
Fontanellen. \gSourcePicture{Gray's Anatomy, Copyright
expired.}} }

% Die \gT{Sutur} ist die endg\"ultig fest verkn\"ocherte Naht
% zwischen den Sch\"adelknochen beim Erwachsenen ${\rightarrow}$ starrer  stabiler 
% Sch\"adel, Schutz  f\"ur Gehirn.

\paragraph[Sch\"adelbasis]{Sch\"adelbasis}~

\gMarried{
% a ... Stirnbein (Os frontale)
% 
% b ... Siebbein (Os ethmoidale)
% 
% c ... Keilbein  (Os sphenoidale)
% 
% d ... Schl\"afenbein(e)  (Os temporale)
% 
% e ... Felsenbein  (Pars petrosa ossis temporalis)
% 
% f ... Hinterhauptsbein  (Os occipitale)
% 
% g ... gro{\ss}es Hinterhauptsloch (Foramen magnum)

Das Gehirn liegt auf der Sch\"adelbasis auf \gZ{(ausser wenn
man auf dem Kopf steht)}. Hirnnerven treten durch die diversen L\"ocher aus dem
Sch\"adel aus. \gEs{Das \gT{gro{\ss}e Hinterhauptsloch}
\gZ{(Foramen magnum)} ist die Durchtrittsstelle des
R\"uckenmarks.} }
{\includegraphics[width=\linewidth]{./images/anatomie/cavernous_sinus.png}
\captionof{figure}{\gSourcePicture{Gray's Anatomy, Copyright
expired.}}}



\gCave{\index{Hirnstammeinklemmung}\gT{Hirnstammeinklemmung}

Das \gEs{gro{\ss}e Hinterhauptsloch} ist
die Durchtrittsstelle des R\"uckenmarks. Bei einer
\gEs{Hirndrucksteigerung} (\gE{Hirn\"odem, Hirnblutung}) ist
es die einzige M\"oglichkeit zur Ausdehnung des Gehirns!
\gEs{Dabei dr\"uckt sich das Gehirn allerdings selber ab},
es kommt zu einer t\"odlichen \gT{Hirnstammeinklemmung!} ${\rightarrow}$ Lebensgefahr! }

 
\begin{minipage}{\linewidth}
\subsection{Die Wirbels\"aule}

\gMarried{
Die Wirbels\"aule besteht aus:

\begin{itemize}
    \item Wirbeln \gZ{(vertebrae)}
    \item \index{Bandscheiben}\gT{Bandscheiben}
    (\gT{Discus}): dienen der "`Sto{\ss}d\"ampfung"'
    \item \gEs{Im Wirbelkanal liegt das R\"uckenmark
    eingebettet.}
\end{itemize}

\vspace{2em}

Sie gliedert sich in:

\begin{itemize}
    \item 7 \gEs{Hals}wirbeln \gZ{(Cervical-Wirbel)}
    \item 12 \gEs{Brust}wirbeln  \gZ{(Thorakal-Wirbel)}
    \item 5 \gEs{Lenden}wirbeln  \gZ{(Lumbal-Wirbel)}
    \item 5 verschmolzene \gEs{Kreuzbein}wirbel \gZ{(Os
    sacrum)}
    \item 3-4 \gEs{Stei{\ss}bein}"`wirbel"'
    \gZ{(Os coccygis)}
\end{itemize}

Die Wirbelsäule hat ein doppelt S-f\"ormig gekr\"ummtes
Achsenskelett, welches den  aufrechten Gang erm\"oglicht.
}
{
\includegraphics[width=0.9\linewidth]{./images/anatomie/wirbelsaeule-lena-edited.png}
\captionof{figure}{\gAuthNotesPicLegalOk{}{}Die Wirbelsäule.
\gSourcePicture{Hirtler.}}}

\end{minipage}





\subsubsection{Die Wirbel: Aufbau und Arten}~

\gMarried{
Prinzipieller Aufbau:

\begin{itemize}
\item Wirbelk\"orper
\item Wirbelbogen  (umschlie{\ss}t das R\"uckenmark)
\begin{itemize}
\item 2 Querforts\"atze
\item 1 Dornfortsatz
\end{itemize}
\end{itemize}


In den jeweiligen Abschnitten sehen die Wirbel
unterschiedlich aus.
}{
\includegraphics[width=\linewidth]{./images/anatomie/Gray94.png}
{\centering\tiny\itshape vorne}
\captionof{figure}{Typischer Wirbelkörper.
\gSourcePicture{Gray's Anatomy, Copyright expired.}}
}





\subsection{Der Brustkorb
(\index{Brustkorb}\index{Thorax}Thorax)}

Der  Brustkorb besteht vorne aus dem \gT{Brustbein}
(\gT{Sternum}) und 12 Rippenpaaren, welche hinten mit der
Wirbels\"aule in gelenkiger Verbindung stehen. 

Die obersten 7 Rippen-Paare werden als 
{\quotedblbase}\gT{echte Rippen}{\textquotedblleft}
bezeichnet. Sie erreichen mit ihren Spitzen das Brustbein und bilden praktisch
{\quotedblbase}Ringe{\textquotedblleft} um die Brusteingeweide. Die 5
unteren Rippenpaare bilden den bogenf\"ormigen unteren Rand des
Brustkorbes. Die Rippenpaare 8-10 sind dabei nur indirekt
\"uber einen Knorpel mit dem Brustbein verbunden
({\quotedblbase}\gT{unechte Rippen}{\textquotedblleft}), das
11. und das 12. Paar endet mit knorpeligen Spitzen frei zwischen den
Bauchmuskeln({\quotedblbase}\gT{fliegende
Rippen}{\textquotedblleft}) bezeichnet werden.

Der Brustkorb enth\"alt und sch\"utzt die Brust\-organe (Lunge, Herz)
sowie die Oberbauchorgane (Leber, Magen). Zwischen den
Rippen spannt sich die \gT{Zwischenrippenmuskulatur}
(wichtig f\"ur \gEs{Atemmechanik}), in diese sind segmentale
Nerven- und Gef\"a{\ss}b\"undel eingelagert. Das Zwerchfell
schließt den Brustraum nach unten hin ab und trennt somit
den Brustraum vom Bauchraum.

\gMarried{Der Thorax besteht aus:

\begin{itemize}
    \item \gEs{Brustbein} (\gT{Sternum})
    \item \gEs{12  Rippenpaaren}
    \begin{itemize}
        \item \gZ{7 "`echte"'}
        \item \gZ{3 "`unechte"'}
        \item \gZ{2 "`fliegende"'}
    \end{itemize}
    \item \gEs{Brustwirbels\"aule} (12 Brust\-wirbel), sie
    steht mit den Rippen in gelenkiger Verbindung
\end{itemize}
(\emph{{\quotedblbase}echt{\textquotedblleft}} bedeutet die
Rippe hat eine direkte Verbindung zum Sternum)}
{\includegraphics[width=\linewidth]{./images/anatomie/thorax.png}

\captionof{figure}{Der knöcherne Thorax.
\gSourcePicture{Mariana Ruiz Villarreal, Public Domain} }}


\subsection{Der \index{Schulterg\"urtel}Schulterg\"urtel
und die obere Extremität}

\begin{figure}[H]
    \begin{center}
    \includegraphics[width=\linewidth]{./images/anatomie/Human_arm_bones_diagram_de-edited.png}
\captionof{figure}{\gAuthNotesPicLegalOk{}{PD}Der
Schultergürtel und die obere Extremität.
\gSourcePicture{Mariana Ruiz Villarreal, Public Domain}}
    \end{center}
\end{figure}

\gParallel{\subsubsection{Der Schultergürtel} 

\ldots besteht aus (von vorne nach
hinten):

\begin{itemize}
    \item \index{Brustbein}\gT{Brustbein}
    (\index{Sternum}\gT{Sternum})
    \item \index{Schl\"usselbein}\gT{Schl\"usselbein}
    \gZ{(Clavicula)}
    \item \index{Schulterblatt}\gT{Schulterblatt} \gZ{(Scapula)}
\end{itemize}

Der Schulterg\"urtel bildet die Gelenkspfanne f\"ur das Schultergelenk,
jeweils auf beiden Seiten.

Die Schulterbl\"atter sind nur durch Muskulatur mit dem Brustkorb
verbunden.
}{
\subsubsection{Die obere Extremität}

\begin{itemize}
    \item \gEs{Schultergelenk}
    \item \gEs{Oberarmknochen}  (\gT{Humerus})
    \item \gEs{Ellenbogengelenk}
    \item \gEs{Unterarmknochen} (2)
    \begin{itemize}
        \item \gEs{Elle}  (\gT{Ulna})
        \item \gEs{Speiche}  (\gT{Radius}), daumenseitig!
    \end{itemize}
    \item Handwurzelgelenk
    \item (8) \gEs{Handwurzelknochen}
    \item (5) \gEs{Mittelhandknochen}
    \item \gEs{Fingergliedknochen} mit 2-3 Gliedern
\end{itemize}
}






\gAuthNotes{GABS: Bild evt. nur mit einer Seite, dafür aber
größer und neben Text?

KOCH: passt so.}


\gFazit{Achtung: Unterscheide: Hand -- Arm!}









\subsection{Der Beckeng\"urtel und die untere Extremität}











\gMarried{
\subsubsection{Der Beckeng\"urtel}


Besteht aus (von hinten nach vorne) :

\begin{itemize}
    \item Kreuzbein (Os sacrum)
    \item Beckenknochen, bestehend aus:
    \begin{itemize}
        \item \gZ{\index{Darmbein}Darmbein}
        \item \gZ{\index{Sitzbein}Sitzbein}
        \item \gZ{\index{Schambein}Schambein}
    \end{itemize}
\end{itemize}
\gZ{Die \index{Symphyse}Symphyse stellt eine straffe
Knorpelverbindung vorne zwischen den beiden Beckenknochen
dar (Dehnungs\-fuge f\"ur Geburt).}

\subsubsection{Die untere Extremit\"at}

\begin{itemize}
    \item \gEs{Hüftgelenk}
    \item \gEs{Oberschenkel}  (\gT{Femur})
    \item \gEs{Kniegelenk}
    \item \gEs{Kniescheibe}
    \item \gEs{Unterschenkelknochen} (2)

    \begin{itemize}
        \item \gEs{Schienbein}  (Tibia) , dick, vorne
        \item \gEs{Wadenbein}  (Fibula), d\"unn, seitlich
        hinten
    \end{itemize}
    \item \gEs{Sprunggelenk}
    \item 8 \gEs{Fu{\ss}wurzelknochen}, \gZ{u.a. Sprungbein,
    Fersenbein}
    \item 5 \gEs{Mittelfu{\ss}knochen}  
    \item \gEs{Zehengliederknochen} mit je 2-3 Gliedern
\end{itemize}
}{
\includegraphics[width=\linewidth]{./images/anatomie/untere-extremitaet.png}
\captionof{figure}{\gAuthNotesPicLegalOk{}{}Der
Beckengürtel und die untere Extremität.
\gSourcePicture{Mariana Ruiz Villarreal, Public Domain}}
}

% TODO STEU: Knöchel erwähnen.
\gA{STEU: Knöchel erwähnen.}

\gFazit{Achtung: Unterscheide Fu{\ss} -- Bein!}

\subsection{Die Muskulatur}

Prinzipiell gibt es 3  Arten von Muskelzellen:

\begin{itemize}
    \item \gT{Quergestreifte (Skelett-)Muskulatur}:
    willk\"urlich steuerbar, trainierbar, er\-m\"ud\-bar
    \item \gT{Glatte  Muskulatur}: Nicht willentlich
    kontrollierbar, unerm\"udlich, ausdauernd
    \item \gT{Herzmuskel}: Sonderstellung zwischen glatt und
    quergestreift, Eigenschaften von beiden: nicht willentlich
    steuerbar, trainierbar, nicht ermüdend, quergestreift.
\end{itemize}

Die \gE{quergestreifte Muskulatur} heisst deswegen
{\quotedblbase}quergestreift{\textquotedblleft}, da man unter dem
Mikroskop tats\"achlich eine Streifung sehen kann, allerdings nicht mit
dem freien Auge. 

\gMarried{Muskel können nur Zug und keinen
Druck ausüben, daher sind \gEs{Gegenspieler} (Agonist u.
Antagonist) erforderlich (z.B. ein Muskel streckt den Arm,
ein \gE{anderer} beugt ihn). Jeder Muskel hat einen Ursprung
und einen Ansatz, diese beiden Punkte bestimmen die Bewegung.
}{
\includegraphics[width=\linewidth]{./images/anatomie/Gray361.png}
\captionof{figure}{\gAuthNotesPicLegalOk{}{}Muskelzug. \gSourcePicture{Gray's
Anatomy, Copyright expired.}}}





%%%%%%%%%%%%%%%%%%%%%%%%%%%%%%%%%%%%%%%%%%%%%%%%%%%%%%%%%%%
%%%%%%%%%%%%%%%%%%%%%%%%%%%%%%%%%%%%%%%%%%%%%%%%%%%%%%%%%%%
%%%%%%%%%%%%%%%%%%%%%%%%%%%%%%%%%%%%%%%%%%%%%%%%%%%%%%%%%%%
%%%%%%%%%%%%%%%%%%%%%%%%%%%%%%%%%%%%%%%%%%%%%%%%%%%%%%%%%%%
\clearpage
\section{Das Herz}

\gMarried{
Das Herz ist ein etwa faustgro{\ss}er \gEs{Hohlmuskel} und
hat die Funktion einer \gEs{Pumpe}. Es liegt \gEs{hinter dem
Sternum\footnote{Brustbein}}\gA{\footnote{\gA{EMHO: statt
Fußzeile: "`Brustbein (Sternum)"'. damit einheitlich.\par
GABS: nein. Hier sollen sie bereits primär den Begriff
"`Sternum"' verwenden udn daran gewöhnt werden.}}} und
teilweise im linken Brustkorb. Unten liegt es auf \gEs{dem Zwerchfell auf}.

Das Herz ist durch eine \gEs{Scheidewand} (Septum) in eine
\gE{rechte} und eine \gE{linke} Seite geteilt.
Umgangssprachlich spricht man dabei jeweils vom
\emph{{\quotedblbase}rechten{\textquotedblleft}} oder
\emph{{\quotedblbase}linken Herz{\textquotedblleft}}. Jede
H\"alfte wird durch \gEs{Klappen} in einen \gT{Vorhof}
(\gT{Atrium}) und eine \gT{Herzkammer} (\gT{Ventrikel})
geteilt.
}{
 \includegraphics[width=\linewidth]{./images/noauth/AASdev20090228-img27.png} 
 \captionof{figure}{\gAuthNotesPicLegalNotOk{img27}{}Das Herz.
 \gSourcePicture{Quelle nicht eruierbar.}}
}



 
\gMarried
{
\paragraph{Herzklappen}

\gAuthNotes{\footnote{\gA{GABS: Wieviel Augenmerk soll auf
die Klappen gelegt werden? KOCH: passt.}}}Insgesamt \gEs{4
Herzklappen} sorgen als Ventile daf\"ur dass das Blut nur in eine Richtung str\"omt. Es gibt 2 Segelklappen und 2 Taschenklappen.
Die Segelklappen befinden sich zwischen den Vorh\"ofen und den
Herzkammern, die Taschenklappen befinden sich an den
Ausfluss\"offnungen der Kammern am \"Ubergang zur Lungenarterie (rechte
Seite) bzw. zur  Aorta (Hauptschlagader, linke Seite). 
}{
 \includegraphics[width=\linewidth]{./images/noauth/AASdev20090228-img28.jpg} 
 \captionof{figure}{\gAuthNotesPicLegalNotOk{img28}{}Das
 Herz.\gAuthNotes{G-FDL !!!, KOCH: Bild gefällt mir nicht. }
\gSourcePicture{Wikipedia,  J\"org Rittmeister. G-FDL, CCAL}}
}


\paragraph{Blutversorgung}

Die Blutversorgung erfolgt \"uber die
\gT{Herzkranzgef\"a{\ss}e} (\gT{Koronargef\"a{\ss}e}) welche
direkt aus der Hauptschlagader abgehen.

\paragraph{Schichten}

3 Schichten (au{\ss}en nach innen):
\gAuthNotes{\footnote{GABS: Schichten relevant? KOCH: ja!}}
\begin{itemize}
\item Epikard (aussen)
\item \gEs{Myokard} (\gEs{Muskel})
\item Endokard (innen)
\end{itemize}



 
% \captionof{figure}{ Schema der
% Herz-Anatomie.  }


\paragraph{Herzfunktion}\label{topic:herzfunktion}

Der Herzmuskel pumpt und entspannt sich beim Erwachsenen
in Ruhe ca. 60 bis 100~mal in der Minute. Bei jeder
Anspannung (Kontraktion) wirft das Herz ca. 70\,ml Blut in die Arterien
aus. Dabei pumpt es ca. 4\,-\,6\,l  Blut (bzw. 1\,x das
gesammte Blutvolumen) in der Minute.

Durch die Pumpfunktion des Herzens ensteht ein
lebenswichtiger Druck im (arteriellen) Gefäßsystem, der
(systolische) \gE{Blutdruck}. Vgl. Kap.
\emph{Vitalfunktionen}, \ref{topic:blutdruck}. 

%
%
\subsection{Reizleitungssystem}

\gAuthNotes{GABS @ STEU: Reizleitungssystem -- wie
ausf\"uhrlich? Pr\"ufungsfrage?

KOCH: Beim Defi schon.

GABS: Defi-Teil noch nciht abgefasst!}

\gMarried{Eigenst\"andiges  Erregungs- und  Reizleitungssystem, streng
hierarchisch, aber jede  Station zur Erregungsbildung  f\"ahig 

bestehend aus:

\begin{itemize}
\item Sinusknoten 
({\sffamily 1}):  
{\quotedblbase}Schrittmacher{\textquotedblleft}, feuert mit einer Eigenfrequenz v. ca. 60/min. elektrische Impulse ab)
\item AV-(AtrioVentricular)-Knoten ({\sffamily 2}): Leitung
vom Vorhof in die Kammern, filtert zu hohe Frequenzen
\item His-B\"undel
\item \gZ{Tawara-}Schenkel (linker/rechter)
\item \gZ{Purkinje-}Fasern
({\quotedblbase}Endstrecke{\textquotedblleft})
\end{itemize}

\gFazit{Der normale Schrittmacher ist der Sinusknoten,
welcher die elektrischen Impulse f\"ur die Vorh\"ofe und die Kammern erzeugt. Man spricht dann
vom Sinusrhythmus.}

Beispiele von typischen Rhythmen werden im Kapitel
Vitalfunktionen vorgestellt \gAuthNotes{Referenz einfügen!}.
}
{
\includegraphics[width=\linewidth]{./images/anatomie/Reizleitungssystem_1.png}

\captionof{figure}{\gAuthNotesPicLegalOk{}{CC-BY}
Reizleitungssystem, Schema. \gSourcePicture{J. Heuser,
basierend auf der Arbeit von Patrick J. Lynch; illustrator;
C. Carl Jaffe; MD; cardiologist Yale University Center for
Advanced Instructional Media. Lizenz: CC-BY}}
}





\gAuthNotes{KOCH: Tabelle mit Normalwerten einfügen

GABS: in Kap Vitalfunktionen}


\section{Der Kreislauf }

\gFazit{\gEs{Herz + Gef\"a{\ss}system + Blut = Kreislauf}}

Der Blutdruck und die Pumpfunktion des Herzens werden im
Kapitel \emph{Vitalfunktionen} unter \ref{topic:blutdruck}
(Blutdruck) sowie im Abschnitt \ref{topic:herzfunktion}
(Herzfunktion) in diesem Kapitel besprochen.

Es wird der \gT{K\"orperkreislauf} (\gE{"`Gro{\ss}er
Kreislauf"'},
"`Hochdrucksystem"') vom \gT{Lungenkreislauf}
(\gE{"`Kleiner Kreislauf"'},
"`Niederdrucksystem"') unterschieden. Das Herz ist f\"ur 
beide Systeme die gemeinsame Pumpe.

Der Kreislauf wird vom \gEs{Stammhirn} gesteuert
(Herzfrequenz, Blutdrugregulation, etc.).

\subsection{Abschnitte des Kreislaufs}
\vspace{-2em}

\begin{gWideParEnv}
\gMarried
{\subsubsection{Kleiner (Lungen-)Kreislauf}
Die obere und untere Hohlvene m\"unden in den \gEs{rechten
Vorhof} des Herzens und liefern sauerstoffarmes Blut aus
dem K\"orper an. Von hier wird das Blut in die \gEs{rechte
Herzkammer} gepumpt, um dann in die \gEs{Lungenarterie} zu
gelangen. In der Lunge ver\"asteln sich die Lungenarterien
zu haard\"unnen Gef\"a{\ss}en, den \gEs{Kapillaren}. Diese
umspannen die Lungenbl\"aschen es findet ein
\gEs{Gasaustausch} statt: Sauerstoff wird aufgenommen und
Kohlendioxid abgegeben. Das jetzt mit Sauerstoff
angereicherte Blut wird \"uber die \gEs{Lungenvenen} in den
\gEs{linken Vorhof} geleitet, der wiederum das Blut in die
\gEs{linke Herzkammer} pumpt.

\vspace{1em}

\subsubsection{Gro{\ss}er Kreislauf}
Das Blut wird aus der linken Herzkammer ausgesto{\ss}en und
gelangt \"uber die die Hauptschlagader
("`\gEs{Aorta}"') in den
K\"orperkreislauf. Aus der Aorta entspringen \gEs{Arterien}
({\quotedblbase}Schlagadern{\textquotedblleft}), die sich
immer weiter verzweigen und an Durchmesser abnehmen. Sie
versorgen den gesamten K\"orper mit Blut. Der durch die
Pumpleistung des Herzens erzeugte Puls ist an oberfl\"achlich
gelegenen Arterien tastbar (Radialis-Puls, Carotis-Puls, ..).
Daher r\"uhrt auch der Name
"`Schlagader"'.

Sie verzweigen sich, weiter an Durchmesser abnehmend,
\"ahnlich dem Ge\"ast eines Baumes, bishin zu
\gEs{Arteriolen} und haard\"unnen Gef\"a{\ss}en, den
Haargef\"a{\ss}en ("`\gEs{Kapillaren}"').
Hier findet wieder ein \gEs{Gasaustausch} statt, diesmal
allerdings spiegelbildlich zum Lungenkreislauf:
\gEs{Sauerstoff wird an das Gewebe abgegeben, im Gegenzug
wird Kohlendioxid abtransportiert}.

Das nun sauerstoffarme Blut gelangt \"uber kleine
\gEs{Venolen} in das Venensystem und flie{\ss}t in den
\gEs{Venen} herzw\"arts. Am Weg zum Herzen nehmen die Venen
an Umfang immer mehr zu, bis sie letztendlich in die gro{\ss}en Venen
({\quotedblbase}\gEs{obere} und \gEs{untere
Hohlvene}{\textquotedblleft}) m\"unden. 
}
{
\includegraphics[width=\linewidth]{./images/noauth/kreislauf-schema.png}

\captionof{figure}{\gAuthNotesPicLegalNotOk{}{Blutkreislauf}Der
Blutkreislauf.}

\gFazit{Rechter Vorhof ${\rightarrow}$ rechte Herzkammer
${\rightarrow}$ Lungenarterie(n) ${\rightarrow}$
Lungenkapillaren ($CO_2$-Abgabe, $O_2$-Aufnahme) ${\rightarrow}$ Lungenvenen ${\rightarrow}$ linker Vorhof 

${\rightarrow}$ linke Herzkammer ${\rightarrow}$ Aorta ${\rightarrow}$
abzweigende Arterien ${\rightarrow}$ Kapillaren
($O_2$-Abgabe, $CO_2$-Aufnahme)  ${\rightarrow}$ Venen
${\rightarrow}$ Hohlvene ${\rightarrow}$ Rechter Vorhof

 ${\rightarrow}$ alles beginnt wieder von vorne
 $\circlearrowright$}}
\end{gWideParEnv}

~


\gAuthNotes{Klappen fehlen in der Beschreibung! KOCH:
passt.}

\subsection{Die Blutgefäße}

Man unterscheidet zwischen:

\begin{description}
    \item[\gT{Arterien}] führen vom Herzen weg, dicke
    Muskelschicht
    \item[\gT{Venen}] führen zum Herzen hin,
    besitzen Venenklappen um Rückfluß\footnote{Rückfluß:
    Fluß entgegen der Flußrichtung} zu vermeiden
    \item[\gT{Kapillargefäße}] d\"unne "`Haargef\"a{\ss}e"'
    in denen der Gasaustausch stattfindet, Verbindung zwischen Arterien
und Venen
\end{description}

\gZ{Arterien besitzen eine dicke Muskelschicht, Venen
hingegen haben eine deutlich d\"unnere Wand als die
muskul\"osen Arterien.} Ein \gT{Shunt} ist eine künstlich
(chirurgisch) angelegte Verbindung zwischen Arterien und
Venen und wird oft bei Dialyse-Patienten eingesetzt
(vgl. {topic:blutdruckmessung}).

%\begin{gWideParEnv}
%\begin{center}
\begin{tabulary}{\linewidth}{lllL}
\toprule
\multicolumn{4}{c}{\gTabHeading{\"Ubersicht wichtiger
Blutgef\"a{\ss}e}}
\\\midrule
\gEs{Hauptschlagader}
 &
\gT{Aorta}
 &
Arterie
 &
Geht von der linken Herzkammer ab, von dort zweigen alle anderen
K\"orperarterien ab
\\\addlinespace
\gEs{Halsschlagader}
 &
\gT{Karotis}
 &
Arterie
 &
1 Paar (2 Stk.), versorgen den Sch\"adel, Puls tastbar seitlich vom
Kehlkopf
\\\addlinespace
\gEs{Herzkranzgef\"a{\ss}e}
 &
\gT{Koronararterien}
 &
Arterie
 &
Versorgen das Herz mit Blut, gehen aus der Aorta ab
\\\addlinespace
\gEs{Lungenarterie}
 &
\gT{Pulmonalarterie}
 &
Arterie 
 &
\gEs{-- $O_2$-arm.} Geht von der rechten Herzkammer weg in
die Lunge zu den Lungenkapillaren 
\\\addlinespace
\gEs{Oberarmarterie}
 &
Brachialarterie
 &
Arterie
 &
Taststelle zum Puls-f\"uhlen beim Baby, Blutdruckmessung
\\\addlinespace
\gEs{Radialarterie}
 &
\gT{A. radialis}
 &
Arterie
 &
Verl\"auft am Unterarm speichenseitig (dort wo der Daumen ist). Puls
f\"uhlbar.
\\\addlinespace
\gEs{Obere Hohlvene}
 &
\gT{Vena cava} sup.
 &
Vene
 &
Obere Vene die in den rechten Vorhof m\"undet
\\\addlinespace
\gEs{Untere Hohlvene}
 &
\gT{Vena cava} inf.
 &
Vene
 &
Untere Vene die in den rechten Vorhof m\"undet. 
\\\addlinespace
\gEs{Lungenvene}
 &
\gT{Pulmonalvene}
 &
Vene 
 &
\gEs{-- $O_2$-reich.} Kommt von den Lungenkapillaren und
m\"undet in den linken Vorhof 
\\\addlinespace
\gEs{Pfortader}
 &
\gZ{Vena portae}
 &
Vene
 &
F\"uhrt O2-armes, aber n\"ahrstoffreiches Blut von Teilen des Darms zur
Leber (zur Entgiftung)
\\\addlinespace
\gEs{Haargef\"a{\ss}e}
 &
\gT{Kapillare}
 &
Kapillare
 &
D\"unne, dicht verzweigte Gef\"a{\ss}e, Verbindung zwischen Arterien
und Venen. Hierfindet der Gas- und N\"ahrstoffaustausch statt. Es gibt
K\"orperkapillaren (K\"orperkreislauf) und Lungenkapillaren
(Lungenkreislauf)
\\\bottomrule
\end{tabulary}
%\end{center}
%\end{gWideParEnv}

\subsection{Das Blut}\index{Blut}\label{topic:blut}

Das Blut macht ca. \gEs{7-8\,\% des K\"orpergewichtes} aus,
das entspricht in etwa \gEs{5-6\,l}.

Es setzt sich zusammen aus:

\begin{itemize}
    \item festen Bestandteilen (Zellen, Blutk\"orperchen)
    \item flüssigen Bestandteilen (\index{Plasma}Plasma)
\end{itemize}

\begin{minipage}{\linewidth}
 \includegraphics[width=\linewidth]{./images/noauth/AASdev20090228-img32.png} 
 \captionof{figure}{\gAuthNotesPicLegalNotOk{img32}{}Die
 Bestandteile des Blutes.
 \gSourcePicture{Quelle nicht eruierbar.}}
\end{minipage}


\paragraph{Blut - Aufgaben}

\begin{itemize}
    \item Sauerstofftransport zu den K\"orperzellen
    \item $CO_2$-Abtransport
    \item Transport von N\"ahrstoffen und Abfallprodukten
    \item W\"armeverteilung
    \item \gZ{pH-Pufferungssystem (S\"aure/Basen-Gleichgewicht)
    (pH-Wert mu{\ss} konstant zwischen 7,35 und 7,45 liegen!)}
\end{itemize}


\subsubsection{Feste (zellul\"are) Bestandteile}

\begin{itemize}
    \item \gT{Rote Blutk\"orperchen} (\gT{Erythrozyten})
    \item \gT{Wei{\ss}e Blutk\"orperchen} (\gT{Leukozyten})
    \item \gT{Blutpl\"attchen} (\gT{Thrombozyten})
\end{itemize}

\paragraph{Erythrozyten}\index{Erythrozyten}

\begin{itemize}
    \item Rote Blutk\"orperchen
    \item Enthalten den roten, eisenhaltigen Blutfarbstoff
    \gT{Hämoglobin}\index{Hämoglobin}, das Hämoglobin ist
    das Transportprotein für Sauerstof
    \item $O_2$-Transport
    \item Tragen Blutgruppenmerkmale (Blutgruppe)
    \item seeeeeehr viele, die h\"aufigsten Blutk\"orperchen
    \gZ{(4.000.000\,/$\mu$m\textsuperscript{3})}
\end{itemize}

\paragraph{Leukozyten}\index{Leukozyten}

\begin{itemize}
    \item Wei{\ss}e Blutk\"orperchen
    \item K\"orperabwehr, Entz\"undungszellen
    \item
    \gZ{(4\,000\,-\,8\,000\,/$\mu$m\textsuperscript{3})}
\end{itemize}

\paragraph{Thrombozyten}\index{Thrombozyten}

\begin{itemize}
    \item \index{Blutpl\"attchen}Blutpl\"attchen
    \item Wichtige Rolle beim Wundverschlu{\ss}
    \item Zelluläre Blutstillung
    \item
    \gZ{(150\,000\,-\,300\,000\,/$\mu$m\textsuperscript{3})}
\end{itemize}

\subsubsection{Die flüssigen Bestandteile --
Das Plasma}\index{Plasma}

\begin{itemize}
    \item 10\% Proteine, Enzyme, Hormone, Elektrolyte, Glucose, Fette
    \item 90\% Wasser
    \item Aufgaben:
    \begin{itemize}
        \item N\"ahrstofftransport
        \item W\"armeverteilung
    \end{itemize}
    \item enth\"alt Proteine und Enzyme f\"ur
    \gEs{Plasmatische Blutgerinnung} (Serum)
    \item Pufferung des pH-Wertes
\end{itemize}

\subsection{Die Blutstillung}

\gMarried{
Die Blutstillung umfasst alle Vorgänge, die zwischen dem Entstehen und dem
Verschluss einer Wunde ablaufen. Sie läuft in zwei Phasen ab: vorläufiger
Wundverschluss durch Bildung eines Thrombozytenpfropfs und endgültiger
Wundverschluss durch Blutgerinnung. 
}{
\includegraphics[width=\linewidth]{./images/anatomie/Bleeding_finger.jpg}
\captionof{figure}{\gAuthNotesPicLegalOk{}{CC-BY}\gSourcePicture{Crystal
(Crystl) from Bloomington, USA
[\url{http://www.flickr.com/people/crystalflickr/}]
Lizenz: CC-BY}} }


\paragraph{Bildung eines Thrombozytenpfropfs} Zunächst
lagern sich an die verletzte und defekte Stelle des Blutgefäßes Thrombozyten an
und verkleben miteinander. Die verklebenden
Blutplättchen bilden einen Thrombozytenpfropf.

\gFazit{Die Bildung des Thrombozytenpropf wird von den zelluären Bestandteilen
des Blutes erledigt.}

\paragraph{Blutgerinnung} \gZ{Die Blutgerinnung ist der wichtigste
Prozess bei der Blutstillung. Sie wird entweder durch die
Gewebeverletzung selber oder durch ein eigenes System von Gerinnungsfaktoren
ausgelöst. Der Vorgang der Blutgerinnung hat zum Ziel,
das im Blutplasma gelöste \gZ{Fibrinogen} zu aktivieren und in ein einen gallertigen
klebrigen und später festen Zustand zu überführen. Um dieses Ziel zu erreichen, steht dem Organismus ein
Gerinnungsystem bestehend aus im Plasma gelösten Gerinnungsfaktoren zur
Verfügung, das kaskadenartig aufgebaut ist. Dieses System funktioniert wie eine
Kette aufgestellter Dominosteine, die nacheinander umfallen, wenn der erste Stein
angestoßen wird. }

\gFazit{Die Blutgerinnung kommt durch im Plasma gelöste Gerinnungsfaktoren
zustande.}

Neben dem System der Blutgerinnung
gibt es auch ein System zur Auflösung von Thromben, die Fibrinolyse. Beide
Systeme müssen sich die Waage halten.

\gFazit{An der Blutstillung sind sowohl die festen Bestandteile des Blutes als
auch das Plasma beteiligt.}

\clearpage
\section{Das Nervensystem}

\paragraph{Unterteilung}~

Man kann das Nervensystem nach seinem Aufbau und nach seiner Funktion
unterteilen.

Vom Blickwinkel des \underline{Aufbaus} wird das Nervensystem grob
unterteilt in ein


\begin{itemize}
\item \gT{Zentrales Nervensystem} (\gT{ZNS}) und ein 
\item \gT{Peripheres Nervensystem}.
\end{itemize}

Das ZNS ist im wesentlichen f\"ur Steuer-, Regel- und Denkfunktionen
zust\"andig sowie Teilweise f\"ur die Weiterleitung von Impulsen und
Informationen. Das periphere Nervensystem ist fast ausschliesslich nur
f\"ur die Weiterleitung von Informationen und Impulsen zu den
Endorganen (oder umgekehrt, von den Endorganen zum ZNS) zust\"andig. 

Betrachtet man die \underline{Funktion} kann man das
periphere Nervensystem in ein

\begin{itemize}
\item \gT{Motorisches} (willk\"urliche Bewegung),
\item \gT{Sensorisches} ("`fühlend"', Reizweiterleitung von
den Sinnesorganen) und ein
\item \gT{Vegetatives} (autonomes, bzw. unwillkürliches)
Nervensystem
\end{itemize}

unterteilen. 

\subsection{Das Zentralnervensystem}

\gMarried{

}{
\includegraphics[width=\linewidth]{./images/anatomie/Brain_human_sagittal_section.png}
\captionof{figure}{Das Hirn im seitlichen Querschnitt.
\gSourcePicture{Patrick J. Lynch, medical illustrator,
\url{http://patricklynch.net}; C. Carl Jaffe, MD,
cardiologist. Lizenz: CC-BY 2.5}} }


\gMarried{Das Zentralnervensystem (ZNS) besteht aus:

\begin{itemize}
    \item Gehirn
    \begin{itemize}
        \item Gro{\ss}hirn
        \item Kleinhirn
        \item Hirnstamm
    \end{itemize}
    \item R\"uckenmark
\end{itemize}

% TODO Hirnteile in grau+ Balken
\gA{STEU: Hirnteile, Balken in grau}

Aufgaben:
\begin{itemize}
    \item Aufnahme  und Verarbeitung  von \"au{\ss}eren und internen  Reizen
    (Schmerz, F\"uhlen, Sehen, H\"oren, {\dots})
    \item Bewusstsein und {\quotedblbase}Bewusstmachung{\textquotedblleft}
    \item Abgabe  von zweckentsprechenden Impulsen (Steuerung, Reaktion
    etc.)
\end{itemize}
}

{\captionof{figure}{\gAuthNotes{Bild fehlt: Hirn
Querschnitt ZNS Beschriftung}}}

\paragraph{ZNS -- Aufgaben}~

\index{Gro{\ss}hirn}\index{Zwischenhirn}\index{Kleinhirn}\index{Hirnstamm}
\index{verl\"angertes R\"uckenmark}\index{R\"uckenmark}

\gT{Gro{\ss}hirn}

\begin{itemize}
    \item Bewu{\ss}tes Denken
    \item Sitz des Bewusstseins
    \item Gefühle
    \item Gedächtnis
\end{itemize}

% \gT{Zwischenhirn}
% \begin{itemize}
% \item Schaltinstanz und Filter R\"uckenmark - Gro{\ss}hirn
% \end{itemize}

\gT{Kleinhirn}

\begin{itemize}
    \item Bewegungskoordination ({\quotedblbase}Makros{\textquotedblleft})
\end{itemize}

\gT{Hirnstamm} und \gT{verl\"angertes R\"uckenmark}

\begin{itemize}
    \item Atem-, Kreislaufzentrum
\end{itemize}

\gT{R\"uckenmark}

\begin{itemize}
    \item Weiterleitung von Impulsen (Leitungsbahnen)
    \item Vom R\"uckenmark gehen die peripheren Nerven aus
    \item (Muskeleigen-)Reflexe
\end{itemize}


\gMarried{\paragraph{Der Schädel besteht aus Schichten}~
%\index{Epiduralraum}
%\index{Dura mater}
\index{Harte Hirnhaut}
%\index{Subduralraum}
%\index{Arachnoidea}
%\index{Subarachnoidalraum}
\gAuthNotes{GABS: Wie sehr hervorgehoben?}
\begin{itemize}
    \item Kopfschwarte
    \item Knochen
%    \item Epiduralraum
    \item Harte Hirnhaut \gZ{(Dura mater)}
%    \item Subduralraum
    \item Arachnoidea (Spinngewebshaut)
%    \item Subarachnoidalraum
%    \item eingebettete Gef\"a{\ss}e
    \item Weiche Hirnhaut \gZ{(Pia mater)}
\end{itemize}}
{
\includegraphics[width=\linewidth]{./images/anatomie/Gray1196-edited.png}
\captionof{figure}{\gAuthNotesPicLegalOk{}{PD-CE}Querschnitt 
die Schädeldecke. \gSourcePicture{Gray's Anatomy, Copyright
expired.}}}


% Aufbau ZNS
% 
% \begin{itemize}
% \item Gro{\ss}hirn (Telencephalon)
% 
% \end{itemize}
% \begin{itemize}
% \item \index{Graue Substanz}Graue Substanz (Rinde)
% 
% \item \index{Wei{\ss}e Substanz}Wei{\ss}e Substanz (Mark)
% 
% \item Balken (Verbindung li/re)
% 
% \end{itemize}
% \begin{itemize}
% \item Zwischenhirn  (Diencephalon)
% 
% \end{itemize}
% \begin{itemize}
% \item Thalamus, Hypothalamus (Hormone, Temperaturregulation etc.)
% 
% \item Hypophyse (Hormone)
% 
% \end{itemize}
% \begin{itemize}
% \item Hirnstamm (Truncus encephali)
% 
% \item Mittelhirn (Mesencephalon)
% 
% \item Br\"ucke (Pons)
% 
% \item verl\"angertes Mark (Medulla oblongata)
% 
% \item Ventrikelsystem
% 
% \end{itemize}
% \begin{itemize}
% \item Liquorgef\"ullte Hohlr\"aume (Ventrikel)
% 
% \item haben nichts mit den Vnetrikel im Herz zu tun!
% 
% \end{itemize}
% Liquor = Gehirnwasser, in Ventrikeln gebildet

Das Hirn und das R\"uckenmark schwimmt in einer
Flüssigkeit, dem \gT{Liquor}. Die
Schädelknochen und der Liquor bieten eine Schutzfunktion. 
Wenn Liquor austritt (z.B. durch Nase, Ohr, ...)  ist das
ein Zeichen f\"ur eine schwerwiegende Verletzung!

\gFazit{Tritt Liquor in die Umwelt aus ist das ein Zeichen
f\"ur eine schwerwiegende Verletzung.} 

\gAuthNotes{KOCH: doppelt? GABS: ja. Schlüsselsätze sollen
doppelt (normales Format + Fazit-Box) vorkommen}

% \paragraph{\gZ{Blut-Hirn-Schranke}} 
% 
% \gZ{Die Blutgef\"a{\ss}wand ist im Hirn besonders dicht. Sie
% filtert daher besonders gut und nur f\"ur bestimmte Molek\"ule ($O_2$,
% $CO_2$, Wasser, \ldots) leicht passierbar. Sie bietet dem Hirn
% Schutz vor sch\"adlichen und giftigen Substanzen. 
% 
% Die Blut-Hirn-Schranke bietet dem Hirn besonderen
% Schutz vor sch\"adlichen Substanzen im Blut.} 

\subsection{PNS -- Peripheres Nervensystem}

\gMarried
{\begin{itemize}
    \item Kommunikation ZNS ${\longleftrightarrow}$
    periphere  Organe
    \item Hirnnerven %(12)
    \item Spinalnerven (gehen vom Rückenmark segmental ab)
    \item elektrische Impulse, chemische Signale (Synapsen)
\end{itemize}

\paragraph{Arten von Nerven}~

\begin{itemize}
    \item motorisch
    \item sensorisch
    \item vegetativ
    \begin{itemize}
        \item \gT{sympathisch}
        \item \gT{parasympathisch}
    \end{itemize}
\end{itemize}}
{\includegraphics[width=\linewidth]{./images/noauth/AASdev20090228-img34.png}
\captionof{figure}{\gAuthNotesPicLegalNotOk{img34}{}R\"uckenmark
mit abgehenden peripheren Nerven. Beachte dass die Nerven
geordnet nach R\"uckenmarkssegmenten abgehen.
\gSourcePicture{Quelle nicht eruierbar.}}

}


\subsubsection{Vegetatives (autonomes)  Nervensystem}
\paragraph{Sympathikus  und Parasymphatikus  sind 
Gegenspieler}~

\gParallel{
\gT{Sympathikus}

Aktiviert den Aufmerksamkeits- und Kampfmodus.

\emph{"`Fight or flight"'} -- \emph{"`Kampf und Flucht"'}
}{
\gT{Parasympathikus} 

Aktiviert den Ruhemodus.

\emph{"`Rest and digest"'} - \emph{"`Rast und Verdauung"'}
}

z.B.: RR + HF

\gParallel{
Sympathikus:  HF ${\uparrow}$, RR ${\uparrow}$
}{
Parasymphatikus: HF ${\downarrow}$, RR ${\downarrow}$
}




\begin{figure}[H]
    \begin{center}
    \includegraphics[width=0.4\linewidth]{./images/noauth/AASdev20090228-img35.png}
    \includegraphics[width=0.4\linewidth]{./images/noauth/AASdev20090228-img36.png} 
    \captionof{figure}{\gAuthNotesPicLegalNotOk{img35 
img36}{}Fight or flight. \gSourcePicture{Quelle nicht
eruierbar.}}
    \end{center}
\end{figure}

\begin{figure}[H]
    \begin{center}
   
    \includegraphics[width=0.8\linewidth]{./images/noauth/AASdev20090228-img37.png} 
    \captionof{figure}{\gAuthNotesPicLegalNotOk{img37}{}
    Nur zur Illustration: Das vegetative Nervensystem hat
    Einflu{\ss} auf viele Organe. Diese Bild
    \underline{nicht lernen}. Es dient nur der
    Verdeutlichung wie wichtig und einflussreich das vegetative
    Nervensystem im K\"orper ist. \gSourcePicture{Quelle nicht eruierbar.}} 
    \end{center}
\end{figure}
 
 
 
 
 
 
 
\clearpage
\section{Die Haut}


Gr\"o{\ss}tes Sinnesorgan und Grenzfl\"ache zur
Au{\ss}enwelt. \gZ{Sie macht ca. \nicefrac{1}{6} des
K\"orpergewichtsaus und ihre Oberfläche beträgt beim
Erwachsenen ca. 1,6 m{\texttwosuperior}.}


\paragraph{Aufgaben}~

\begin{itemize}
    \item Schutz vor \"au{\ss}eren Einfl\"ussen (mechanisch, Erreger,
    UV-Strahlung etc.)
    \item Schutz vor Wasserverlust/Austrocknung
    \item Schutz vor W\"armeverlust / 
    \item Temperaturregulation (Schwei{\ss})
    \item Sekretionsorgan (Schwei{\ss}, Talg)
    \item Sinnesorgan (Tastsinn, Temperatursinn)
\end{itemize}

\paragraph{Die Haut ist in Schichten aufgebaut}~
 
\gAuthNotes{Haut: Bild Querschnitt: muss vereinfacht werden,
lat. Ausdr\"ucke weg bzw. grau}

\index{Oberhaut}\index{Lederhaut}\index{Unterhaut}

\gMarried{
\begin{itemize}
    \item \gT{Oberhaut} \gZ{(Epidermis)}
    \begin{itemize}
        \item Hornschicht, geschichtetes Plattenepithel
    \end{itemize}
    \item \gT{Lederhaut} \gZ{(Corium, Dermis)}
    \begin{itemize}
        \item Hautanhangsgebilde
        \begin{itemize}
            \item Schwei{\ss}dr\"usen
            \item Talgdr\"usen
            \item Haarfollikel, Nagelmatrix
        \end{itemize}
        \item Tastk\"orperchen, freie Nervenendigungen (Schmerz)
        \item \gZ{Kollagenfasern (St\"utznetzwerk)}
        \item Arteriolen, Venolen (Versorgung der Haut, Temperaturregulation)
    \end{itemize}
    \item \gT{Unterhaut} \gZ{(Subcutis)}
    \begin{itemize}
        \item Depotfett, W\"armeisolierung
    \end{itemize}
\end{itemize}
}{
\hspace{-0.5cm}\includegraphics[width=1.2\linewidth]{./images/noauth/AASdev20090228-img38.png}
\captionof{figure}{\gAuthNotesPicLegalNotOk{img38}{Haut
Querschnitt}Die Haut im Querschnitt.
\gSourcePicture{Quelle nicht eruierbar.}} }



\paragraph{Die Haut ist wichtig f\"ur die
\gEs{Thermoregulation}}~

Neben der Lunge dient vor allen Dingen  die Haut als wichtigstes Organ
zur \gEs{Thermoregulation}: Die Blutgef\"a{\ss}e  der  Haut
sind stark ineinander verwickelt und k\"onnen sich bei Bedarf stark erweitern. Die
stark durchblutete  Haut strahlt somit W\"arme ab. Weiters wird durch
die Verdunstungsk\"alte  des Schwei{\ss}es  ebenfalls W\"arme
abgegeben. 

\paragraph{Wundliegen}
\index{Wundliegen}\index{Dekubitus}\label{topic:dekubitus}Wenn
der Mensch nicht fähig ist in ausreichendem Umfang regelmäßig seine Lage zu ändern kann es zum wundliegen kommen. Die Wunden nennt man dann
\gT{Dekubitus}. \gZ{Besonders gefährdet sind jene Stellen,
an denen ein Knochen knapp unter der Haut liegt
(Durckstelle). Betroffen sind vor allem alte oder
pflegebedürftige bettlägrige Menschen, bewusstlose
Patienten bzw. sedierte Intensivpatienten sowie Patienten
mit Querschnittsläsionen.}


\clearpage
\section{Die Atemwege}

Die Atemwege gliedern sich in: 
\gAuthNotes{KOCH: mag Nasenrachenraum und Mundrachenraum
streichen, evt. wie Bild?}

\gMarried
{
\gT{Obere Atemwege}:

\begin{itemize}
    \item \gEs{Nasen}- und \gEs{Mundhöhle}
    \item \gEs{Rachen} (gemeinsamer Weg von Nahrung und
    Atem!)
    \item \index{Kehlkopf}\gEs{Kehlkopf} (\gT{Larynx}) mit
    \gEs{Kehldeckel}
\end{itemize}

}
{
\gT{Untere Atemwege}:
\begin{itemize}
    \item \index{Luftr\"ohre}\gEs{Luftr\"ohre}
    (\index{Trachea}\gT{Trachea})
    \item \index{Hauptbronchus}\gT{Hauptbronchus} (li./re.)
    \item \index{Bronchioli}\gT{Bronchioli}
    \item \index{Lungenbl\"aschen}\gEs{Lungenbl\"aschen}
    (\index{Alveolen}\gT{Alveolen})
\end{itemize}
}
%{\gTempPicInWork{Atemwege Übersicht img39}}

\begin{figure}[h]
    \begin{center}
    \includegraphics[width=\linewidth]{./images/anatomie/Respiratory_system_complete_de-edited.png} \captionof{figure}{\gAuthNotesPicLegalOk{}{PD}Die Atemwege, Übersicht.
\gSourcePicture{Mariana Ruiz Villarreal, Public domain}}
    \end{center}
\label{fig:respiraktionstrakt-uebersicht}
\end{figure}



%\begin{multicols}{2}
\paragraph{Obere Atemwege}

Der Weg der Einatemluft beginnt in der Nase bzw. im Mund. Die Nase ist
mit Schleimhaut ausgekleidet, kn\"ocherne Vorspr\"unge sorgen f\"ur
eine Vorw\"armung der Luft. In der Schleimhaut be\-finden sich die
Sinneszellen f\"ur den Geruchssinn. 

\gZ{In den Sch\"adelknochen befinden sich luft\-ge\-f\"ullte
Hohlr\"aume, die so\-genannten \gT{Neben\-h\"ohlen} (Stirn-,
Kiefer\-h\"ohlen, Siebbeinzellen). Diese stehen in Verbindung mit dem Nasenrachenraum.
Entz\"undet sich die Schleimhaut spricht man von der
\gE{Neben\-h\"ohlen\-ent\-z\"undung}, welche oft langwierig
und schmerzhaft sein kann.} 

Die Luft gelangt nun in den \gT{Mund-Rachen\-raum} und den
\gT{Rachen}. \gEs{Hier kreuzt der Weg der Nahrung mit dem
Weg der Atemluft!} Die Nahrung gelangt vom Mundraum in die
hinten gelegene Speiser\"ohre, die Luft m\"ochte vom Nasen\-rachen\-raum in die vorne gelegene Luftr\"ohre. 

Um ein Verschlucken von Nahrung in die Lunge zu verhindern
verschlie{\ss}t der am Kehlkopf gelegene \gT{Kehldeckel}
w\"ahrend des Schluckaktes die Luftr\"ohre. Dieser Reflex ist lebenswichtig! 

Im \gE{Kehlkopf}  befinden sich die \gT{Stimmb\"ander},
welche f\"ur die Lautbildung beim Sprechen wichtig sind.

\paragraph{Untere Atemwege}

In der \gT{Luftr\"ohre} (\gT{Trachea}) wird die vom Rachen
und Kehlkopf kommende Luft zur Lunge weitergeleitet. Die Luftr\"ohre teilt sich ein einen
\gE{rechten} und einen \gE{linken} \gT{Hauptbronchus},
welche zu den jeweiligen Lungenfl\"ugeln ziehen. Jeder Hauptbronchus verzweigt sich immer weiter
bis schlu{\ss}endlich sehr d\"unne \gT{Bronchioli} in die
\gT{Lungenbl\"aschen} (\gT{Alveolen}) m\"unden. Hier findet
der Gasaustausch statt.

Die Abschnitte der Atemwege, in denen kein Gasaustausch
stattfindet und der nur dem Transport dient, bezeichnet man
als {\quotedblbase}\index{Totraum}\gT{Totraum}{\textquotedblleft}.

%\end{multicols}


\subsection{Die Lunge ist f\"ur einen Teil des
Gasaustausches zust\"andig}

Die Lunge hat 2 Fl\"ugel, welche sich jeweils in 3 (rechts) bzw. 2
(links) Lappen gliedern: 

\begin{center}
\begin{tabularx}{\linewidth}[H]{XX}
\toprule
\multicolumn{2}{c}{Die Lunge}
\\\toprule
\multicolumn{1}{c}{Rechter Fl\"ugel}
 &
\multicolumn{1}{c}{Linker Fl\"ugel}
\\
\begin{itemize}
    \item \gEs{3} Lappen
    \begin{itemize}
        \item Oberlappen
        \item Mittellappen
        \item Unterlappen
    \end{itemize}
\end{itemize}
&
\begin{itemize}
    \item \gEs{2} Lappen
    \begin{itemize}
        \item Oberlappen
        \item Unterlappen
    \end{itemize}
\end{itemize}
\\\bottomrule
\end{tabularx}
\end{center}

Die Oberfl\"ache ist mit dem \gT{Lungenfell}, der
\gE{lungenseitigen} \gT{Pleura}, \"uberzogen. In einem
Lungefl\"ugel verzweigt sich der Hauptbronchus in immer weitere kleinere Bronchien bis schlie{\ss}lich die Bronchioli in
die \gT{Lungenbl\"aschen} (\gT{Alveolen}) enden. 

\subsubsection{Die Alveolen sind die Endst\"ucke
des Atemwegs und sind f\"ur den Gasaustausch zust\"andig}

\emph{Siehe Abb. \ref{fig:respiraktionstrakt-uebersicht}
rechts oben.}

In den \gT{Lungenbl\"aschen} (\gT{Alveolen}) findet der
\index{Gasaustausch}\gE{Gasaustausch} zwischen der
Einatemluft und dem Blut statt. Sauerstoff gelangt dabei aus den Alveolen in das Blut und
Kohlendioxid wird im Gegenzug vom Blut an die Alveolen
abgegeben.

Dabei besteht zwischen Luft und Blut kein direkter Kontakt, die Gase
m\"ussen Hindernisse wie die Alveolenwand und die Blutgef\"a{\ss}wand
\"uberwinden. \gE{Vergr\"o{\ss}ert sich dieser Weg aufgrund
einer Erkrankung ist die Sauerstoffaufnahme vermindert!} 

\gFazit{
Sauerstoff ($O_2$):   Alveolen ${\rightarrow}$ Blut

Kohlendioxid ($CO_2$):  Blut ${\rightarrow}$ Alveolen
}


\subsubsection{Die Pleura}\label{topic:pleura}

\gMarried{
\begin{itemize}
    \item \gE{Thoraxseitige} Pleura (\gT{Rippenfell}, fest mit
    Brustwand und Zwerchfell verwachsen, kleidet damit die Brusth\"ohle aus)
    \item \gE{Lungenseitige} Pleura (\gT{Lungenfell},
    \"uberzieht die Lungenoberfl\"ache)
    \item Aufgrund eines Unterdrucks haften die beiden Pleuren aneinander,
    sind aber gegeneinander verschieblich (wie 2 Glasplatten)
    \item Geringe Menge Pleuraflüssigkeit ${\rightarrow}$ Adhäsion der
    beiden Blätter!
    \item Normalerweise keine Verwachsungen der beiden Bl\"atter
\end{itemize}
}
{\includegraphics[width=\linewidth]{./images/noauth/AASdev20090228-img41.png}
\captionof{figure}{\gAuthNotesPicLegalNotOk{img41}{Pleura}Pleura.
\gSourcePicture{Quelle nicht eruierbar.}}}



\paragraph{Welche Funktion hat aber nun die Pleura?}

Da die eine Seite mit dem Brustkorb, die andere Seite mit der Lunge
verwachsen ist und der Raum zwischen den beiden Seiten luftdicht
abgechlossen ist haften die beiden Seiten aneinander, aber sind
gegeneinander verschieblich. (Es gibt ja keine
Verwachsungen der beiden Seiten und die Pleuraflüssigkeit fungiert sogar ein
bisschen als "`Schmiermittel"'.) 

Wenn sich nun der Brustkorb oder das Zwerchfell bei der
Atmung bewegt, bewegt sich natürlich das Rippenfell mit (es ist ja mit dem Brustkorb
und dem Zwerchfell verwachsen). Da der Pleurazwischenraum
luftdicht ist, muss sich zwangsläufig aber auch das Lungenfell (an dem
wiederum die Lunge angewachsen ist) mitbewegen. Und wenn sich das Lungenfell
mitbewegt wird automatisch auch die Lunge wie eine Ziehharmonika
"`aufgezogen"' und Luft
"`angesaugt"'. 

\subparagraph{Frage:} \emph{Was würde passieren wenn der
Zwischenraum zwischen den Pleurablättern aufgrund einer Verletzung (Z.B. ein Loch
nach einer Stichwunde in die Brust) nicht mehr luftdicht
ist?} Antwort: s. Fußzeile (\footnote{Antwort: Die Lunge
würde in sich zusammenfallen, da sie ja nicht mehr aufgespannt wird. Siehe
"`Pneumothorax"'.}) 



\subsection{Atemmechanik}
\label{topic:atemmechanik}

Wenn man sich den Thorax annähernd als Zylinder vorstellt,
so kann eine Volumsvergrößerung (Einatmung) durch eine
Vegrößerung des Wuerschnittes (Heben der
Zwischenrippenmuskulatur) oder Vergrößerung der Höhe
(Anspannung des Zwerchfells) erfolgen.

\begin{equation}
V = A \times h
\end{equation}
\begin{equation}
\mbox{Volumen} = \mbox{Fläche} \times \mbox{Höhe}
\end{equation}

\gAuthNotes{KOCH: entweder wird das genauer erklärt wie die
Funktion der Pleura oder es wird ganz einfach mündlich
erklärt.

GABS: Sollte schon schriftlich erklärt werden \ldots wir
wollen ja auch immer dass sie die "`Atemhilfsmuskulatur"'
etc. unterstützen, dann sollten sie wenigstens auch
nachlesen können "`wie man richtig atment"'

KOCH: soll erklärt werden.}

\includegraphics[width=\linewidth]{./images/noauth/AASdev20090228-img42.png}
\captionof{figure}{\gAuthNotesPicLegalNotOk{img42}{Atemmechanik}Atemmechanik.
\gSourcePicture{Quelle nicht eruierbar.}}

Für die Atemmachanik ist ein gutes Zusammenspiel von Skelett, Muskeln, Gewebe
und Pleura(spalt) wichtig. 

\gFazit{Ein bis zum Hals Verschütteter erstickt.}

\paragraph{Atemmuskulatur} Der wichtigste Atemmuskel ist das \gEs{Zwerchfell}.
In Ruhe erledigt es fast die gesammte Atemarbeit. Bei verstärkter Atmung kommen
die Zwischenrippenmuskeln zum Einsatz \cite{LippertAnatomie7}.

\paragraph{Atemhilfsmuskulatur} Neben den Hauptatemmuskeln (Zwerchfell,
Zwischenrippenmuskulatur) gibt es noch eine Reihe weiterer Atemhilfsmuskeln die
die Atmung unterstützen. Einige haben nur dann eine Funktion wenn der
Schultergürtel fixiert ist. Dies erreicht man zum Beispiel durch aufstützen der Arme. Vgl. dazu \emph{'Asthma',
\ref{topic:asthma-bronchiale}} und \emph{'COPD', \ref{topic:copd}}.

\paragraph{Pleura} Die Funktion der Pleura ist unter \ref{topic:pleura} erklärt.

\gZ{Thoraxwand und Lunge besitzen elastische Eigenschaften: Sie werden gedehnt
und entwickeln elastische Rückstellkräfte. An der Lunge können zwei Mechanismen ein
Zusammenfallen des Lungengewebes bewirken. Erstens sind im Lungenzwischengewebe
reichlich elastische Fasern eingelagert, die die Eigenelastizität der Lunge
bewirken. Durch die Koppelung von Lunge und
Thoraxwand am Pleuraspalt entsteht ein elastisches Gesamtsystem. 

Bei entspannter
Atemmuskulatur und offenem Kontakt mit der Umgebungsluft wird sich ein
Gleichgewicht der elastischen Kräfte einstellen; der Brustkorb befindet sich in
der Atemruhelage. Durch den Zug der Lunge, die sich verkleinern möchte, und den
Gegenzug der Thoraxwand entsteht ein ständiger Unterdruck im Pleuraspalt.}

\subparagraph{Atemzyklus} 

\gZ{Volumenschwankungen des Brustraums bewirken
gleichzeitig auch Volumenschwankungen im Inneren der Lunge. Eine Vergrößerung des Thoraxinnenraums
wird durch Kontraktion der Atemmuskulatur bewirkt. Dieser Vorgang wird als aktive
Phase der Atmung bezeichnet. Dabei kommt es vor allem zu einer Abflachung und
einem Tiefertreten des Zwerchfells, aber auch zu einer Anhebung der Rippen durch
die äußere Zwischenrippenmuskeln. Als Folge entsteht in der Lunge im Vergleich
zur Umgebungsluft ein Unterdruck, der durch Einströmen von Luft ausgeglichen
wird. Thoraxwand und Lunge werden dabei gedehnt. 

Am Ende einer ruhigen Einatmung
erschlafft die Einatmungsmuskulatur. Die elastischen Rückstellkräfte führen das
Thoraxvolumen wieder in die Ausgangslage zurück. Dabei kommt es zu einer
Druckerhöhung in der Lunge, die eine Luftbewegung hin zur Umgebungsluft bewirkt.
Diese erfolgt in aller Regel passiv, d. h. ohne aktive
muskuläre Anstrengung. Dieser Vorgang erfolgt in Ruhe bei einem Erwachsenen etwa zwölfmal in der Minute
und wird als Atemfrequenz bezeichnet.}


\gFazit{Der wichtigste Atemmuskel ist das \gEs{Zwerchfell}!}

\subsection{\gZ{Atemfrequenz, -tiefe und -volumen}}

Die Atemfrequenz, die Atemtiefe und das Atemvolumen
werden im Kapitel Vitalfunktionen unter
\gSee{topic:atemvolumina} besprochen.

\subsection{Luft}

Die Luft, welche wir atmen, ist ein Gasgemisch aus
Stickstoff, Sauerstoff, Kohlendioxid und Edelgasen.
Zwischen der normalen Einatem- und Ausatemluft besteht
besonders hinsichtlich des Sauerstoff- und
Kohlendioxid-Anteils ein Unterschied, siehe
Tab. \ref{tab:luft-gas}.

\begin{center}
\begin{tabular}[H]{ccc}
\toprule
\centering Raumluft & & Ausatemluft
\\\midrule
\gEs{21\,\%}  & Sauerstoff ($O_2$) & \gEs{16\,\%}
\\
79\,\% & Stickstoff ($N$) & 79\,\%
\\
\gEs{0,03\,\%} & Kohlendioxid ($CO_2$) & \gEs{4\,\%} 
\\
1\,\% & Edelgase  & 1\,\%
\\\bottomrule
\end{tabular}
\captionof{table}{Gasgehalt in der Raum- und
Ausatemluft.\label{tab:luft-gas}}
\end{center}


\clearpage
\section{Der Verdaungstrakt -- Von der Schüssel in die
Schüssel}
\index{Magen-Darm-Trakt}\index{Gastrointestinaltrakt}

\emph{Syn.: \gT{Magen-Darm-Trakt},
\gT{Gastrointestinaltrakt} (\gT{GIT})}

\begin{itemize}
    \item Die Verdauung dient dem Stoffwechsel.
    \item Proze{\ss}:
    \begin{itemize}
        \item Nahrungsaufnahme
        \item Weiterverarbeitung im Magen-Darm-Trakt
        \item Transport der N\"ahrstoffe und Energietr\"ager  ins Blut
        \item Zellaufbau und Abbau von Abfallstoffen
    \end{itemize}
    \item \"Uber die Verdauung gewinnt der K\"orper  also einerseits die
    n\"otigen {\quotedblbase}Bausteine{\textquotedblleft},  andererseits
    auch die n\"otige Energie um  Zellen aufbauen und in Funktion erhalten
    zu k\"onnen.
\end{itemize}

\paragraph{\"Ubersicht}~

\begin{figure}[H]
    \includegraphics[width=0.8\linewidth]{./images/anatomie/Digestive_system_diagram_de.png}
    \captionof{figure}{\gAuthNotesPicLegalOk{}{}Schema des Verdauungstraktes.
    \gSourcePicture{Mariana Ruiz Villarreal, Public domain}}
\end{figure}

\begin{minipage}{\linewidth}
\captionof{table}{Übersicht über den Verdauungstrakt --
\emph{"`Von der Schüssel in die Schüssel."'}\label{tab:verdauungstrakt-uebersicht}}
\begin{tabulary}{\linewidth}[H]{LLL}
\toprule
\multirow{3}{8em}{\gEs{Mundhöle}}
 &
 Zunge
 &
\multirow{2}{15em}{Zerkleinerung der Nahrung und Vermischung
mit Speichel}
\\\hhline{~-~}
 &
Z\"ahne 
&
\\\hhline{~-~}
 &
Speicheldrüsen
&
\\\hhline{---}
Rachen
 &
 &
Kreuzung des Speisebreis mit dem Luftweg

Kehldeckel verschliesst beim Schlucken die Luftr\"ohre
\\\hhline{---}
\gEs{Speiser\"ohre} (\gZ{\"Osophagus})
 &
 &
Speisebrei wird mittels Muskelbewegung zum Magen befördert
\\\hhline{---}
\gEs{Magen}
 &
 &
Vermischung mit Magensäure

Magens\"aure zerlegt Nahrung

Pf\"ortner gibt Portionsweise Speisebrei an den Zw\"olffingerdarm ab

\\\hhline{---}
\multirow{4}{8em}{\gEs{D\"unndarm}}
&
\gT{Zw\"olffingerdarm} \gZ{Duodenum}
 &
Gallenfl\"ussigkeit (${\leftarrow}$Leber) und Bauchspeichel
(${\leftarrow}$Pankreas) wird zugesetzt
\\\hhline{~--}
 &
\gT{Krummdarm} \gZ{Jejunum}
 &
Spaltung von N\"ahrstoffen (Eiwei{\ss}, Kohlehydrate, Fett) und
N\"ahrstoffaufnahme
\\\hhline{~-~}
 &
\gT{Leerdarm} \gZ{Ileum}
 &
\\\hhline{---}
\multirow{6}{8em}{\gEs{Dickdarm}\\(\gT{Colon})}
&
\gEs{Blinddarm}
&
Sackgasse
\\
&
\gEs{Wurmfortsatz} (\gT{Appendix}) 
 &
Anh\"angsel der Sackgasse, kann sich entz\"unden.
\\\hhline{~--}
 &
\gEs{Aufsteigender} Teil
 &
\multirow{4}{\columnwidth}{Wasserentzug, \\Eindickung, \\
Bakterienbesiedlung, \\Produktion und Speicherung des
Stuhls} 
\\\hhline{~-~}
 &
\gEs{Querender} Teil
 &
\\\hhline{~-~}
 &
\gEs{Absteigender} Teil
 &
\\\hhline{~-~}
 &
\gEs{S-f\"ormig} gebogener Teil (\gZ{Sigma})
 &
\\\midrule
\gEs{Mastdarm} (\gT{Rektum})
 &
 &
{\quotedblbase}Speicherung{\textquotedblleft} des ausscheidungsbereiten
Stuhls
\\\hhline{---}
\gEs{Anus}
 &
 &
Absetzen des Stuhles

\\\bottomrule
\end{tabulary}

\end{minipage}


\subsection{Mundhöhle}

\begin{description}
    \item[Zähne] zerkleinern die Nahrung
    \item[Speicheldrüsen] produzieren Speichel um Nahrung
    gleitfähig zu machen und Nahrung mit Verdauungsenzymen \gZ{(für Zucker und Stärke)} zu versetzen.
    \item[Zunge] schiebt Nahrung in Schlund weiter,
    verschlie{\ss}t beim Schlucken Kehldeckel (sonst  \gEs{Aspiration!})
\end{description}



\subsection{Die Speiser\"ohre (Ösophagus)}

\gMarried
{
\begin{itemize}
    \item ca. 24cm langer, dreischichtiger Muskelschlauch
    \gZ{(oberes Drittel quergestreifte Muskulatur, untere
    \nicefrac{2}{3} glatt)}
    \item Verbindung Mund $\rightarrow$ Magen
    \item "`schiebt"' Nahrung durch
    \gEs{Muskelkontraktionen} (\gT{Peristaltik})
    weiter
    \item hier findet keine Verdauung statt
\end{itemize}
}
{
\centering
\includegraphics[width=0.3\linewidth]{./images/noauth/AASdev20090228-img44.png} 
\includegraphics[width=0.3\linewidth]{./images/noauth/AASdev20090228-img45.png}
\captionof{figure}{\gAuthNotesPicLegalNotOk{img44
img45}{}Ösophagus. \gSourcePicture{Quelle nicht eruierbar.}}
}




\subsection{Der Magen}

\gZ{(Gaster, Ventriculus)}

\gAuthNotes{KOCH: Fließtext!}
\begin{itemize}
    \item 2-3\,l Magensaft
    \item Schleim (sch\"utzt  Magenschleimhaut)
    \item Salzs\"aure (Bakterienfalle, pH 1,0)
    \item \gZ{Pepsin  (Enzym f\"ur Eiwei{\ss}-Verdauung)}
    \item Wandmuskulatur \gEs{durchmischt} den Speisebrei
    mit der Magens\"aure
    \item portionsweise Abgabe von Speisebrei durch den
    Magenpf\"ortner \gZ{(Pylorus)} in den Zw\"olffingerdarm
\end{itemize}


\gParallel{
\centering\includegraphics[width=0.7\linewidth]
{./images/noauth/AASdev20090228-img46.png} 
\captionof{figure}{\gAuthNotesPicLegalNotOk{img46
img45}{}Lage des Magens. \gSourcePicture{Quelle nicht
eruierbar.}}
}{
 \includegraphics[width=\linewidth]{./images/noauth/AASdev20090228-img47.png} 
 \captionof{figure}{\gAuthNotesPicLegalNotOk{img47}{}Der
 Magen, der Pf\"ortner und das erste St\"uck
 des Zw\"olffingerdarms im Durchschnitt.
 \gSourcePicture{Quelle nicht eruierbar.}}
}




\subsection{Zw\"olffingerdarm (Duodenum)}

\begin{itemize}
\item Beginn des D\"unndarmes, ca 25\,cm lang
\item Verdauungsenzyme aus Leber und Pankreas f\"ur Fett, Eiwei{\ss},
Zucker m\"unden hier ein
\end{itemize}

\gParallel{
\includegraphics[width=\linewidth]{./images/noauth/AASdev20090228-img48.png}
\captionof{figure}{\gAuthNotesPicLegalNotOk{}{}Die Lage des
Zwöffingerdarms im Körper.} }{
\includegraphics[width=\linewidth]{./images/noauth/AASdev20090228-img49.png}
\captionof{figure}{\gAuthNotesPicLegalNotOk{}{}Der
Zwölffingerdarm und die Bauchspeicheldrüse.} }

\gAuthNotes{KOCH: sind beides wirklich Enzyme?

GAB: Ja.}


\subsection{Die \index{Leber}Leber} 

\gAuthNotes{KOCH: fehlt: was ist sie, eine Drüse, ein
Hohlmuskel, \ldots } 

\gZ{("`Hepar"')}

\gMarried
{\begin{itemize}
    \item Die Leber liegt im rechten Oberbauch und hat einen linken und
    rechten Lappen.
    \item Abbau von Giftstoffen (die Leber nimmst durch die
    Giftstoffe mit der Zeit selbst auch Schaden
    (Leberentzündung, Leberzirrhose)
    \item "`Chemiefabrik"' des Körpers: Produktion von vom
    Körper benötigten Eiweißen
    \item {\quotedblbase}Verdauungsdr\"use{\textquotedblleft}
    \begin{itemize}
        \item produziert \index{Gallenfl\"ussigkeit}Gallenfl\"ussigkeit
        ({\quotedblbase}Galle{\textquotedblleft}), welche in der Gallenblase
        gespeichert wird: Fettverdauung
        \item Von der Leber gehen die Galleng\"ange, an den Galleng\"angen
        h\"angt die
    \end{itemize}
    \item \index{Gallenblase}Gallenblase: Nur Speicher f\"ur
    Gallefl\"ussigkeit
    \item speichert Fett \& Zucker
    \item erh\"alt \"uber \index{Pfortader}\gT{Pfortader}
    (Vena portae) n\"ahrstoffreiches Blut vom Darm zwecks Entgiftung
\end{itemize}}
{
\includegraphics[width=\linewidth]{./images/noauth/AASdev20090228-img50.png}
}

\gAuthNotes{GABS: Erklärung/Bild Pfortader einfügen? fragt
STEU}




\gDias{\includegraphics[width=\linewidth]{./images/noauth/AASdev20090228-img51.png}
\captionof{figure}{\gAuthNotesPicLegalNotOk{img51}{}Leber
Ansicht von vorne.
\gSourcePicture{Quelle nicht eruierbar.}}}
{\includegraphics[width=\linewidth]{./images/noauth/AASdev20090228-img52.png}
\captionof{figure}{\gAuthNotesPicLegalNotOk{img52}{}Leber,
Ansicht von unten mit Gallenblase.
\gSourcePicture{Quelle nicht eruierbar.}}}
{\includegraphics[width=\linewidth]{./images/noauth/AASdev20090228-img53.png}
\captionof{figure}{\gAuthNotesPicLegalNotOk{img53}{}Gallenblase.
\gSourcePicture{Quelle nicht eruierbar.}}}


 

\subsection{Bauchspeicheldr\"use}\index{Bauchspeicheldr\"use}

(\index{Pankreas}\gT{Pankreas})

Das Pankreas ist ca.\,15-22\,cm lang und liegt hinter dem
Magen. Es hat einen \gEs{Ausf\"uhrungsgang}, welcher
zusammen mit dem Gallengang \gEs{in den Zw\"olffingerdarm}
(Duodenum) mündet.

\paragraph{Aufgaben} 

Das Pankreas hat \gEs{zwei} wesentliche Funktionen:

\begin{description}
    \item[Enzymproduktion] f\"ur Verdauung
    ${\rightarrow}$ Zw\"offingerdarm
    \item[Hormonproduktion] Ausschüttung ins Blut
    \begin{itemize}
        \item \index{Insulin}\gT{Insulin} ${\rightarrow}$ Blutzucker
        ${\downarrow}$
        \item \gZ{Glukagon ${\rightarrow}$ Blutzucker ${\uparrow}$}
    \end{itemize}
\end{description}

%\gTempPicInWork{Gallensystem}

\subsection{Krumm- und Leerdarm -- Jejunum und Ileum}
\index{Krummdarm}\index{Leerdarm}

\gT{Jejunum}, \gT{Ileum}

\gMarried
{\begin{itemize}
    \item ca. 3\,m lang
    \item Oberfl\"achenvergr\"o{\ss}erung durch Darm\-zotten
    \item \gEs{N\"ahrstoffaufnahme}
    \item \gZ{beweglich, am Darmgekr\"ose (Mesenterium)
    aufgeh\"angt}
\end{itemize}}
{
\includegraphics[width=0.6\linewidth]{./images/noauth/AASdev20090228-img56.png}
\captionof{figure}{\gAuthNotesPicLegalNotOk{img56}{}Dünndarm
Querschnitt.} }





\subsection{Dickdarm (Colon)}

Der Dickdarm (Colon) ist ca. 1,5\,m lang und wird in
verschiedene Abschnitte unterteilt: 

\begin{center}
\begin{tabulary}{\linewidth}[H]{llL}
\toprule
\multirow{6}{8em}{\gEs{Dickdarm}\\(\gT{Colon})}
&
\gEs{Blinddarm}
&
Sackgasse
\\
&
\gEs{Wurmfortsatz} (\gT{Appendix}) 
 &
Anh\"angsel der Sackgasse, kann sich entz\"unden.
\\\hhline{~--}
 &
\gEs{Aufsteigender} Teil
 &
\multirow{4}{\columnwidth}{Wasserentzug, Eindickung, \\
Bakterienbesiedlung, \\Speicherung} 
\\\hhline{~-~}
 &
\gEs{Querender} Teil
 &
\\\hhline{~-~}
 &
\gEs{Absteigender} Teil
 &
\\\hhline{~-~}
 &
\gEs{S-f\"ormig} gebogener Teil (\gT{Sigma})
 &
\\\bottomrule
\end{tabulary}
\end{center}

\gMarried{
Der Leerdarm m\"undet in den Dickdarm. Diese M\"undung ist wie eine
\gE{T-Kreuzung} gebaut: \gE{Nach oben} geht es in den
\gEs{aufsteigenden Teil} des Dickdarmes, \gE{nach unten} in
den \gEs{Blinddarm}. Wie der Name des Blinddarmes schon sagt
endet er blind, er ist daher eine Sackgasse. Am Blinddarm
angeh\"angt findet man den \gEs{Wurmfortsatz}
(\gT{Appendix}) welcher sich h\"aufig entz\"undet
(\gT{Appendiz\underline{itis}\footnote{Das Anhängsel
\emph{"`-itis"' zeigt immer eine Entzündung an, vgl.
Kapitel Hygiene}}}). Der Wurmfortsatz erf\"ullt beim
Menschen keine nennenswerte Funktion, ausser der
Arbeitsplatzsicherung von An\"asthesisten und Chirurgen
(wegen der vielen Blinddarm-Operationen).
}{\includegraphics[width=0.8\linewidth]{./images/noauth/AASdev20090228-img57.png}

\captionof{figure}{\gAuthNotesPicLegalNotOk{img57}{}Lage
des Dickdarms und des Appendix. \gSourcePicture{Quelle
nicht eruierbar.}
}
}





\clearpage
\section{Sonstige Bauchorgane}

\subsubsection{\index{Bauchfell}Bauchfell}

\begin{itemize}
    \item Die gesamte Bauchh\"ohle ist mit  Bauchfell
    (\index{Peritoneum}Peritoneum) ausgekleidet
    \item d\"unner \"Uberzug
    \item {\quotedblbase}Aufh\"angung{\textquotedblleft} der Bauchorgane
    \item zus\"atzlich gro{\ss}es und kleines Netz
    \item Kann sich Entz\"unden (\gT{Bauchfellentz\"undung}
    = \gT{Peritonitis}): zB bei Darmverletzung,
    Bauchfelldialyse
\end{itemize}


\gMarried
{Der Bauch "`aufgeklappt"': Nach oben geklappt sieht man das
 "`Gro{\ss}e Netz"', auf die Seiten
und nach unten die Bauchmuskulatur. 

Die hellen, fleischfarbenen Schlingen sind D\"unndarmschlingen, die
dicken braunen Schlingen sind der aufsteigende und der quere Teil des
Dickdarmes.}
{
\includegraphics[width=\linewidth]{./images/noauth/AASdev20090228-img58.png}
\captionof{figure}{\gAuthNotesPicLegalNotOk{}{}Der Bauch, geöffnet, mit aufgeklapptem
Bauchfell.}}


\subsection{Die Milz} 

\gZ{Syn: Lien, Splen.}



\gMarried{
\begin{itemize}
\item hat nichts  mit der Verdauung  zu tun 
\item 12\,cm, oval, von \gEs{Kapsel} umgeben
\item im linken Oberbauch unter den Rippen gelegen
\item Teil des Immunsystems
\item entfernt alte Erythrozyten, beim Embryo auch Blut\-bildung
\item sehr gut durchblutet, Blutspeicher (Gefahr:
\index{Milzri{\ss}}Milzri{\ss})
\end{itemize}
}{
\centering\includegraphics[width=0.6\linewidth]{./images/noauth/AASdev20090228-img60.png}
\captionof{figure}{\gAuthNotesPicLegalNotOk{img60}{}Milz:
Lage.
\gSourcePicture{Quelle nicht eruierbar.}}
}



\gCave{\underline{Achtung}: sog. {\quotedblbase}\index{Zweizeitiger
Milzriss}\gT{Zweizeitiger Milzriss}{\textquotedblleft}:

Zuerst reisst die Milz ein, die Kapsel bleib erhalten, es blutet zuerst in die Kapsel und der Druck steigt.
Nach einer gewissen Zeit reisst auch die Kapsel und es blutet nun
ungehindert aus der inneren Wunde. \gEs{Der Patient wirkt
zuerst {\textquotedblleft}ganz normal{\textquotedblleft}, unerwartet und
pl\"otzlich verf\"allt er.}} 



\subsection{Die Niere(n) (Ren)}

\gMarried{
\begin{itemize}
    \item paarig angelegtes Organ
    \item stark durchblutet, empfindlich auf RR- Abfall ${\rightarrow}$
    Funktion ${\downarrow}$
    \item reguliert
    \begin{itemize}
        \item \gEs{Wasser-} und \gEs{Elektrolythaushalt}
        \item Volumshaushalt
        \item pH (S\"aure/Basen)
    \end{itemize}
    \item Geht im Alter, besonders bei Diabetikern, kaputt ${\rightarrow}$
    {\quotedblbase}\gT{Chronische
    Niereninsuffizienz}{\textquotedblleft} ${\rightarrow}$
    {\quotedblbase}\gT{Blutw\"asche}{\textquotedblleft}
    (\index{Dialyse}\gT{Dialyse})
\end{itemize}
}{
\centering\includegraphics[width=0.8\linewidth]{./images/noauth/AASdev20090228-img61.png}
\captionof{figure}{\gAuthNotesPicLegalNotOk{img61}{}Niere.
\gSourcePicture{Quelle nicht eruierbar.}}
}



\paragraph{Aufbau}~

\begin{itemize}
\item Rinde, Mark
\item Hilus (Stiel)
\item Gef\"a{\ss}e
\item Nierenbecken mit Anschluss an den Harntrakt
\end{itemize}

\gAuthNotes{Niere: Anatomie und Funktion muss noch
ausgearbeitet werden.}


\gAuthNotes{KOCH: grau:}
\paragraph{\gZ{Mikroskopischer Aufbau}}

\gZ{
\begin{itemize}
    \item Nephron ist die funktionielle EInheit (ca. 1
    Mio/Niere)
    \begin{itemize}
        \item Glomerulus: Hier wird das Blut gefiltert und Plasmawasser
        abgepresst (ca. 150\,l\,/\,Tag)
        \item Schleifenapparat (Tubulus): Das abgepresste Plasma\-wasser wandert
        den Schleifen\-apparat entlang, hier werden die meisten wichtigen
        Stoffe und der Gro{\ss}\-teil des Wassers wiederum zur\"uckgeholt,
        sodass am Ende nur mehr 1,5 -- 2~l Harn pro Tag in der Harnblasen
        landen.
    \end{itemize}
\end{itemize}
}

\gZ{
\paragraph{Funktionsweise}

\begin{itemize}
    \item 1500 l Blut pro Tag wird
    {\quotedblbase}bearbeitet{\textquotedblleft}
    \item Glomerula: Filtration ${\rightarrow}$ 150\,l
    \index{Prim\"arharn}Prim\"arharn (enth\"alt Elektrolyte, etc.) wird
    gewonnen und fliesst durch den Schleifenapparat
    \item Schleifenapparat: {\quotedblbase}Recycling{\textquotedblleft},
    Gutes kommt zur\"uck ins Blut, der Rest fliesst weiter Richtung
    Harnblase ${\rightarrow}$ 1,5-2\,l (End-)Harn/Tag
\end{itemize}
}

\subsection{Die Nebenniere(n) (Adren)}

% \gTempPicMissing{Nebennieren}

\gZ{Paariges, jeweils am oberen Nierenpol gelegenes Organ.

Hat nichts  mit dem Harntrakt oder der eigentlichen Nierenfunktion zu
tun.

\gEs{Hormonproduktion}
\gZ{
\begin{itemize}
    \item Cortison
    \item Sexualhormone (Androgene)
    \item Adrenalin , Noradrenalin
\end{itemize}}

Die Nebennieren sitzen wie {\quotedblbase}M\"utzen{\textquotedblleft}
auf den Nieren, haben mit deren Funktion allerdings
\gE{nichts} zu tun.
}

\subsection{Harnableitendes System -- Harntrakt}

\gMarried
{Der Harntrakt besteht aus 
\begin{itemize}
    \item Nieren inkl. Nierenbecken
    \item \index{Harnleiter}je 1 Harnleiter  \gZ{(Ureter)},
    30\,cm lang, m\"undet  in die
    \item \index{Blase}Blase 
    \item Harnr\"ohre  (\gZ{Urethra,} Mann: 25\,cm, Frau
    4\,cm)
\end{itemize}

Die Harnr\"ohre verl\"auft beim Mann von der Harnblase durch die
Prostata und den Penis und m\"undet an der Eichel in die Aussenwelt.
Bei einer Verg\"o{\ss}erung der Prostata kann der Harnabflu{\ss} bis
zur Unpassierbarkeit erschwert sein.  

Bei der Frau ist die Harnr\"ohre k\"urzer als beim Mann und
m\"undet oberhalb des Scheideneinganges. 

\captionof{figure}{\gAuthNotesPicLegalOk{}{Lena}Harntrakt
mit Nieren (rechte Niere im Querschnitt mit Nierenbecken),
Harnleiter, Blase und Harnr\"ohre. \label{fig:harntrakt}
\gSourcePicture{Lena Hirtler}}
}
{~

\includegraphics[width=\linewidth]{./images/anatomie/anatomie-harntrakt-edited.png}
} 





\subsection{Was gibt es noch?}

Ausser den "`klassischen Bauchorganen"' liegen auch andere
Organssysteme im Bauchraum auf die nicht vergessen werden
sollte. Besonders betrifft das die weiblichen
Fortpflanzungsorgane wie Eierstöcke, Eileiter und Uterus.
Diese werden allerdings im Kapitel "`Geschlechtsorgane"'
gesondert besprochen. 

\clearpage
\section{Geschlechtsorgane}
\subsection{M\"annliche Geschlechtsorgane}

\begin{itemize}
    \item innen:
    \begin{itemize}
        \item Hoden
        \begin{itemize}
            \item Spermienproduktion
            \item Geschlechtshormon-Produktion
        \end{itemize}
        \item Nebenhoden: Spermienreifung und Lagerung
        \item Samenstrang, besteht aus
        \begin{itemize}
            \item Samenleiter
            \item Nerven
            \item Gef\"a{\ss}e
            \item durchzieht Leistenkanal
        \end{itemize}
        \item \index{Prostata}\gT{Prostata}
        (Vorsteherdr\"use)
        \begin{itemize}
            \item produziert Samenfl\"ussigkeit (Beweglichkeit der Spermien)
            \item \gEs{Vergr\"o{\ss}ert sich im Alter} und
            kann die \gEs{Harnr\"ohre abdr\"ucken}, es kommt
            zum Harnstau ${\rightarrow}$ h\"aufger Notfall!
        \end{itemize}
    \end{itemize}
    \item au{\ss}en:
    \begin{itemize}
        \item Glied (Penis): Schwellk\"orper, Harnsamenr\"ohre, Begattungsorgan
        \item Hodensack \gZ{(Skrotum)}: umh\"ullt Hoden und
        Nebenhoden
        \begin{itemize}
            \item Hodenstiel kann sich verdrehen
            ${\rightarrow}$ \gCa{Notfall!}
        \end{itemize}
    \end{itemize}
\end{itemize}

\begin{figure}[H]
    \begin{center}
     \includegraphics[width=7cm]{./images/anatomie/anatomie-genitale-maennlich-edited.png}
    \captionof{figure}{\gAuthNotesPicLegalOk{}{}Die
    männlichen Geschlechtsorgane im
    Querschnitt. \gSourcePicture{Lena Hirtler}}
    \end{center}
\end{figure}
 

\subsection{Weibliche Geschlechtsorgane}

\begin{itemize}
    \item innen:
    \begin{itemize}
        \item \gEs{Eierst\"ocke} \gZ{(Ovar)}: Reifung der
        Eizellen
        \item \gEs{Eileiter} \gZ{(Tuben)}: Transport des
        reifen Eis in Richtung Geb\"armutter
        \item \gEs{Geb\"armutter}
        (\index{Uterus}\gT{Uterus}): dehnbarer Hohlmuskel,
        mit Schleimhaut ausgekleidet,
        \begin{itemize}
            \item Schleimhaut wird regelm\"a{\ss}ig aufgebaut und wieder
            abgesto{\ss}en (Menstruation, Regelblutung)
            \item Hier nistet sich die befruchtete Eizelle
            ein.
            \item Der Gebärmutterhals (\gT{Cervix}) führt
            in Richtung der Scheide.
            \item Die M\"undung der Geb\"armutter in die
            Scheide ist der \gT{Muttermund}.
        \end{itemize}
        \item \gEs{Scheide} (\gT{Vagina})
    \end{itemize}
    \item au{\ss}en:
    \begin{itemize}
        \item Schamlippen \gZ{(Labien)}:
        go{\ss}e/\"au{\ss}ere, kleine/innere
        \item \gT{Klitoris}: entspricht der m\"annlichen
        Eichel
    \end{itemize}
\end{itemize}


\begin{figure}[H]
    \begin{center}
    \includegraphics[width=7cm]{./images/anatomie/anatomie-genitale-weiblich-edited.png}
    \captionof{figure}{\gAuthNotesPicLegalOk{}{Lena
    }Weibliche Geschlechtsorgane, Querschnitt.
    \gSourcePicture{Lena Hirtler}}
    \end{center}
\end{figure}

 

\subsubsection{Der weibliche Zyklus}\label{topic:weiblicher-zyklus}

\begin{itemize}
    \item Der weibliche Zyklus dauert 28 Tage
    \item gerechnet von Menstruation zu Menstruation  (Monatsblutung )
    \item Uterusschleimhaut wird aufgebaut und wieder abgesto{\ss}en
    \item hormongesteuert
    \begin{itemize}
        \item \"Ostrogen, Progesteron
    \end{itemize}
    \item ca. am 14. Tag: Eisprung
    \item Befruchtung nein:
    \begin{itemize}
        \item Zyklus l\"auft normal weiter, Menstruation
    \end{itemize}
    \item Befruchtung ja:
    \begin{itemize}
        \item Eizelle nistet sich im Uterus ein, Zyklus gestoppt
    \end{itemize}
\end{itemize}

Was kann dabei schief gehen? Wenn sich die Eizelle schon im Ei\-leiter
einnistet kommt es zu einer
\index{Eileiter\-schwanger\-schaft}\gEs{Eileiter\-schwanger\-schaft}
die lebens\-bedrohlich ist (Eileiter rei{\ss}t und die Frau verblutet).

\begin{minipage}{\linewidth}
\includegraphics[width=\linewidth]{./images/noauth/AASdev20090228-img65.png}
\captionof{figure}{\gAuthNotesPicLegalNotOk{img65}{}Der weibliche Zyklus. Oben:
Follikelreifung. Unten: Auf- und Abbau der Uterusschleimhaut. Am Tag 0 des Zyklus wird die Uterusschleimheut abgestossen, es kommt zur
Regelblutung, danach wird die Schleimhaut wieder aufgebaut. Am Tag 14
erfolgt der Eispung. Kommt es zu keiner Befruchtung wird die
Schleimheut erneut am Tag 28 abgestossen. Nach einer
Befruchtung wandert die Eizelle weiter in den Uterus und nistet sich dort ein.
\gSourcePicture{Quelle nicht eruierbar.}}
\end{minipage}



